% !TEX root = print.tex
\chapter{Daubechies wavelets}
\label{ch:daubechies}

This chapter constructs the compactly supported orthonormal wavelets of
Daubechies~\cite{Daubechies} on $\R$ and describes the Besov spaces $\Besov{s}{p}{q}(\R)$ by
their coefficients. A mask of order $N$ is a real polynomial $L$ of degree less than $N$ with
$L(1) = 1$ and $|L(e^{-2\pi i \xi})|^2 = P_N(\sin^2 \pi\xi)$, where
$P_N(y) = \sum_{k < N} \binom{N-1+k}{k} y^k$. The scaling function $\varphi$ is the inverse
Fourier transform of the infinite product of the symbol of the mask. Everything is proved with
the Fourier transform of integrable functions: the refinement equation, the compact support
(through the cascade algorithm), the orthonormality of the integer translates, and the
reproduction of the polynomials of degree less than $N$. The regularity of $\varphi$ comes from
the transfer operator of Eirola~\cite{Eirola} and Villemoes~\cite{Villemoes}, which gives
$\varphi \in H^s$ for $s < 1$ if $N = 2$ and for $s < 3 - \log_2 3$ if $N = 3$, and, for the
first wavelet $N = 2$, from the matrices of Daubechies and Lagarias~\cite{DaubechiesLagarias},
which give the H\"older exponent $2 - \log_2(1 + \sqrt 3)$. With Jackson's and Bernstein's
inequalities the Daubechies coefficients then characterise $\Besov{s}{p}{q}(\R)$.

\section{The Daubechies polynomial}

This section proves the B\'ezout identity of the Daubechies polynomial $P_N$ and the bounds on
$P_N$ that the later sections use. At $y = \sin^2 \pi\xi$ the identity becomes the trigonometric
identity behind the quadrature mirror filter identity of the masks. Throughout,
$\binom{2N-1}{N}$ stands for $\binom{-1}{0} = 1$ if $N = 0$.

The polynomial $P_N$ consists of the first $N$ terms of the binomial series
$(1 - y)^{-N} = \sum_{k \geq 0} \binom{N-1+k}{k} y^k$.

\begin{definition}[The Daubechies polynomial]
\label{def:daub-poly}
\lean{Daubechies.daubPoly, Daubechies.eval_daubPoly, Daubechies.daubPoly_zero,
  Daubechies.daubPoly_one, Daubechies.daubPoly_two, Daubechies.daubPoly_three,
  Daubechies.natDegree_daubPoly_le, Daubechies.eval_daubPoly_succ_eq_one_add}
\leanok
For $N \in \N$ the \emph{Daubechies polynomial} of order $N$ is the real polynomial
\[
  P_N(y) = \sum_{k=0}^{N-1} \binom{N-1+k}{k} y^k
\]
of degree at most $N - 1$. Thus $P_0 = 0$, $P_1 = 1$, $P_2(y) = 1 + 2y$ and
$P_3(y) = 1 + 3y + 6y^2$, and for every $M \in \N$ the constant term of $P_{M+1}$ is $1$:
\[
  P_{M+1}(y) = 1 + \sum_{k=0}^{M-1} \binom{M+k+1}{k+1} y^{k+1} .
\]
\end{definition}

\begin{lemma}[Values at $0$ and $1$]
\label{lem:daub-poly-values}
\lean{Daubechies.eval_zero_daubPoly, Daubechies.eval_one_daubPoly}
\leanok
\uses{def:daub-poly}
For every $N \geq 1$, $P_N(0) = 1$ and $P_N(1) = \binom{2N-1}{N}$.
\end{lemma}

\begin{proof}
\leanok
Only the constant term $\binom{N-1}{0} = 1$ contributes to $P_N(0)$. For $P_N(1)$, write
$\binom{N-1+k}{k} = \binom{N-1+k}{N-1}$ and apply the hockey-stick identity
$\sum_{i=0}^{n} \binom{i+m}{m} = \binom{n+m+1}{m+1}$ with $n = m = N - 1$.
\end{proof}

The coefficients of $P_N$ and $P_{N+1}$ are linked by Pascal's rule.

\begin{lemma}[The recurrence]
\label{lem:daub-poly-recurrence}
\lean{Daubechies.one_sub_mul_eval_daubPoly_succ, Daubechies.one_sub_X_mul_daubPoly_succ,
  Daubechies.choose_two_mul_eq_two_mul_choose}
\leanok
\uses{def:daub-poly}
For every $N \in \N$ and every $y \in \R$,
\[
  (1 - y) P_{N+1}(y) = P_N(y) + \binom{2N-1}{N} y^N - \binom{2N}{N} y^{N+1} ,
\]
and the same identity holds for polynomials. For every $N \geq 1$,
$\binom{2N}{N} = 2 \binom{2N-1}{N}$.
\end{lemma}

\begin{proof}
\leanok
On the left the constant term is $1$, the coefficient of $y^{k+1}$ for $k < N$ is
$\binom{N+k+1}{k+1} - \binom{N+k}{k} = \binom{N+k}{k+1}$ by Pascal's rule, and the top term is
$-\binom{2N}{N} y^{N+1}$. On the right,
$P_N(y) + \binom{2N-1}{N} y^N = \sum_{k=0}^{N} \binom{N-1+k}{k} y^k$ has the same
coefficients. Pascal's rule $\binom{2N}{N} = \binom{2N-1}{N-1} + \binom{2N-1}{N}$ and the
symmetry $\binom{2N-1}{N-1} = \binom{2N-1}{N}$ give the second identity.
\end{proof}

\begin{theorem}[The B\'ezout identity]
\label{thm:daub-bezout}
\lean{Daubechies.eval_daubPoly_bezout_succ, Daubechies.eval_daubPoly_bezout,
  Daubechies.daubPoly_bezout}
\leanok
\uses{def:daub-poly}
For every $N \geq 1$ and every $y \in \R$,
\[
  (1 - y)^N P_N(y) + y^N P_N(1 - y) = 1 ,
\]
and the same identity holds for polynomials.
\end{theorem}

\begin{proof}
\leanok
\uses{lem:daub-poly-recurrence}
Induction on $N$; for $N = 1$ the identity reads $(1 - y) + y = 1$. By
\Cref{lem:daub-poly-recurrence}, from $N$ to $N + 1$ the first summand changes by
$y^N (1 - y)^N \bigl(\binom{2N-1}{N} - \binom{2N}{N} y\bigr)$ and the second by
$y^N (1 - y)^N \bigl(\binom{2N-1}{N} - \binom{2N}{N} (1 - y)\bigr)$. The changes add up to
$y^N (1 - y)^N \bigl(2 \binom{2N-1}{N} - \binom{2N}{N}\bigr) = 0$. Two real polynomials with the
same value at every real number are equal.
\end{proof}

\begin{lemma}[Bounds for $P_N$]
\label{lem:daub-poly-bounds}
\lean{Daubechies.eval_daubPoly_nonneg, Daubechies.one_le_eval_daubPoly,
  Daubechies.eval_daubPoly_le_eval_one, Daubechies.eval_daubPoly_le}
\leanok
\uses{def:daub-poly}
Let $N \in \N$. For every $y \geq 0$, $P_N(y) \geq 0$, and $P_N(y) \geq 1$ if $N \geq 1$. For
every $y \in [0, 1]$,
\[
  P_N(y) \leq P_N(1) \quad\text{and}\quad P_N(y) \leq 1 + \binom{2N-1}{N} y .
\]
\end{lemma}

\begin{proof}
\leanok
\uses{lem:daub-poly-values}
The coefficients of $P_N$ are nonnegative, and for $N \geq 1$ the constant term is $1$. For
$y \in [0, 1]$, $y^k \leq 1$ for every $k$ and $y^k \leq y$ for $k \geq 1$, and the coefficients
of the terms with $k \geq 1$ add up to $P_N(1) - 1 = \binom{2N-1}{N} - 1$
(\Cref{lem:daub-poly-values}). For $N = 0$, $P_N = 0$.
\end{proof}

\begin{lemma}[A binomial bound]
\label{lem:daub-choose-bound}
\lean{Daubechies.choose_lt_four_pow}
\leanok
For every $N \geq 2$, $\binom{2N-1}{N} < 4^{N-1}$.
\end{lemma}

\begin{proof}
\leanok
Let $N = M + 1$. The coefficients $\binom{2M+1}{i}$ with $i \leq M$ form one half of the row
$2M + 1$ of Pascal's triangle, so $\sum_{i=0}^{M} \binom{2M+1}{i} = 2^{2M} = 4^M$. This sum
contains the term $\binom{2M+1}{M} = \binom{2M+1}{M+1}$ and, since $M \geq 1$, the further term
$\binom{2M+1}{0} = 1$, so $\binom{2N-1}{N} = \binom{2M+1}{M+1} \leq 4^M - 1$.
\end{proof}

Equivalently, $N - \log_4 \bigl(2 \binom{2N-1}{N}\bigr) > \frac12$ for every $N \geq 2$.

\paragraph{The trigonometric form.}
On the circle the Daubechies polynomial is evaluated at $\sin^2 \pi\xi$: a mask $L$ of order $N$
has $|L(e^{-2\pi i \xi})|^2 = P_N(\sin^2 \pi\xi)$.

\begin{definition}[The trigonometric Daubechies polynomial]
\label{def:daub-fun}
\lean{Daubechies.daubFun, Daubechies.daubFun_two, Daubechies.daubFun_three}
\leanok
\uses{def:daub-poly}
For $N \in \N$ and $\xi \in \R$ let $P(\xi) = P_N(\sin^2 \pi\xi)$; the order $N$ is fixed and
omitted from the notation. For $N = 2$, $P(\xi) = 1 + 2 \sin^2 \pi\xi$, and for $N = 3$,
$P(\xi) = 1 + 3 \sin^2 \pi\xi + 6 \sin^4 \pi\xi$.
\end{definition}

\begin{lemma}[Symmetries of $P$]
\label{lem:daub-fun-symmetry}
\lean{Daubechies.daubFun_periodic, Daubechies.daubFun_neg, Daubechies.daubFun_add_half,
  Daubechies.continuous_daubFun}
\leanok
\uses{def:daub-fun}
For every $N \in \N$ the function $P$ is continuous, and for every $\xi \in \R$
\[
  P(\xi + 1) = P(\xi), \qquad P(-\xi) = P(\xi), \qquad
  P\bigl(\xi + \tfrac12\bigr) = P_N(\cos^2 \pi\xi) .
\]
\end{lemma}

\begin{proof}
\leanok
$\sin(\pi(\xi + 1)) = -\sin \pi\xi$, $\sin(-\pi\xi) = -\sin \pi\xi$ and
$\sin(\pi(\xi + \frac12)) = \cos \pi\xi$, and $P(\xi)$ is a polynomial in $\sin \pi\xi$.
\end{proof}

\begin{proposition}[Bounds for $P$]
\label{prop:daub-fun-bounds}
\lean{Daubechies.daubFun_nonneg, Daubechies.daubFun_le_choose, Daubechies.daubFun_le,
  Daubechies.daubFun_zero, Daubechies.one_le_daubFun}
\leanok
\uses{def:daub-fun}
For every $N \in \N$ and every $\xi \in \R$,
\[
  0 \leq P(\xi) \leq \binom{2N-1}{N} \quad\text{and}\quad
  P(\xi) \leq 1 + \binom{2N-1}{N} \pi^2 \xi^2 .
\]
If $N \geq 1$, then moreover $P(0) = 1$ and $P(\xi) \geq 1$ for every $\xi \in \R$.
\end{proposition}

\begin{proof}
\leanok
\uses{lem:daub-poly-bounds, lem:daub-poly-values}
\Cref{lem:daub-poly-bounds} at $y = \sin^2 \pi\xi \in [0, 1]$, with $P_N(1) = \binom{2N-1}{N}$
and $P_N(0) = 1$ (\Cref{lem:daub-poly-values}), and $\sin^2 \pi\xi \leq \pi^2 \xi^2$.
\end{proof}

\begin{theorem}[The trigonometric B\'ezout identity]
\label{thm:daub-fun-bezout}
\lean{Daubechies.daubFun_bezout}
\leanok
\uses{def:daub-fun}
For every $N \geq 1$ and every $\xi \in \R$,
\[
  \cos(\pi\xi)^{2N} P(\xi) + \sin(\pi\xi)^{2N} P\bigl(\xi + \tfrac12\bigr) = 1 .
\]
\end{theorem}

\begin{proof}
\leanok
\uses{thm:daub-bezout, lem:daub-fun-symmetry}
\Cref{thm:daub-bezout} at $y = \sin^2 \pi\xi$, where $1 - y = \cos^2 \pi\xi$ and
$P(\xi + \frac12) = P_N(1 - y)$ (\Cref{lem:daub-fun-symmetry}).
\end{proof}

The symbol $m_0$ of a mask of order $N$ has $|m_0(\xi)|^2 = \cos(\pi\xi)^{2N} P(\xi)$, and the
identity becomes the quadrature mirror filter identity
\[
  |m_0(\xi)|^2 + \bigl|m_0\bigl(\xi + \tfrac12\bigr)\bigr|^2 = 1 .
\]

\section{Masks and filters}

This section defines the Daubechies masks with their filter coefficients and symbols, and proves
the identities between the filter coefficients on which the scaling function and the wavelet
rest: the quadrature mirror filter identity, the orthogonality of the even shifts, the sum
rules, and their counterparts for the dual filter. The first Daubechies mask, D4, is explicit.

\paragraph{The unit circle.}
A symbol is a polynomial evaluated at the point $e^{-2\pi i \xi}$ of the unit circle; the sign
of the exponent is the one of the Fourier transform.

\begin{lemma}[The exponentials $e^{-2\pi i \xi}$]
\label{lem:daub-exp}
\lean{Daubechies.expFreq, Daubechies.norm_expFreq, Daubechies.expFreq_ne_zero,
  Daubechies.expFreq_eq_exp_mul, Daubechies.expFreq_re, Daubechies.expFreq_im,
  Daubechies.expFreq_zero, Daubechies.expFreq_add, Daubechies.expFreq_pow,
  Daubechies.expFreq_intCast, Daubechies.expFreq_one, Daubechies.expFreq_periodic,
  Daubechies.expFreq_half_add_int, Daubechies.conj_expFreq, Daubechies.contDiff_expFreq,
  Daubechies.continuous_expFreq, Daubechies.norm_expFreq_sub_one_le,
  Daubechies.norm_expFreq_pow_sub_one_le}
\leanok
The function $\xi \mapsto e^{-2\pi i \xi}$ on $\R$ is smooth and has period $1$. For all
$\xi, \eta \in \R$, every $k \in \N$ and every $n \in \Z$:
\begin{enumerate}
  \item $|e^{-2\pi i \xi}| = 1$, $\operatorname{Re} e^{-2\pi i \xi} = \cos 2\pi\xi$ and
    $\operatorname{Im} e^{-2\pi i \xi} = -\sin 2\pi\xi$;
  \item $\overline{e^{-2\pi i \xi}} = e^{-2\pi i (-\xi)}$ and
    $e^{-2\pi i (\xi + \eta)} = e^{-2\pi i \xi} e^{-2\pi i \eta}$;
  \item $(e^{-2\pi i \xi})^k = e^{-2\pi i k \xi}$, $e^{-2\pi i n} = 1$ and
    $e^{-2\pi i (1/2 + n)} = -1$;
  \item $|e^{-2\pi i k \xi} - 1| \leq 2\pi k |\xi|$.
\end{enumerate}
\end{lemma}

\begin{proof}
\leanok
The exponential turns sums into products and commutes with conjugation, $e^{-2\pi i n} = 1$
for $n \in \Z$, and $e^{-\pi i} = -1$. The bound is $|e^{ix} - 1| \leq |x|$ for real $x$, at
$x = -2\pi k \xi$.
\end{proof}

\begin{lemma}[The factor $(1 + z)/2$ on the unit circle]
\label{lem:daub-half-factor}
\lean{Daubechies.one_add_expFreq_div_two, Daubechies.norm_one_add_expFreq_div_two,
  Daubechies.normSq_one_add_expFreq_div_two}
\leanok
For every $\xi \in \R$,
\[
  \frac{1 + e^{-2\pi i \xi}}{2} = e^{-\pi i \xi} \cos \pi\xi ;
\]
in particular $|1 + e^{-2\pi i \xi}|/2 = |\cos \pi\xi|$ and
$|1 + e^{-2\pi i \xi}|^2/4 = \cos^2 \pi\xi$.
\end{lemma}

\begin{proof}
\leanok
\uses{lem:daub-exp}
$2 \cos \pi\xi = e^{\pi i \xi} + e^{-\pi i \xi}$; multiply by $e^{-\pi i \xi}$, which has
modulus $1$ (\Cref{lem:daub-exp}).
\end{proof}

\paragraph{Masks.}
The factor $((1 + z)/2)^N$ of the mask polynomial produces the zeros of the symbol at the
half-integers and the sum rules of the filter coefficients.

\begin{definition}[Daubechies masks]
\label{def:daub-mask}
\lean{Daubechies.Mask, Daubechies.Mask.one_le}
\leanok
\uses{def:daub-fun}
Let $N \in \N$. A \emph{Daubechies mask of order $N$} is a real polynomial $L$ of degree less
than $N$ with $L(1) = 1$ and
\[
  |L(e^{-2\pi i \xi})|^2 = P(\xi) = P_N(\sin^2 \pi\xi) \quad\text{for every } \xi \in \R .
\]
Every Daubechies mask has order $N \geq 1$.
\end{definition}

\begin{definition}[Mask polynomial, filter coefficients and symbol]
\label{def:daub-symbol}
\lean{Daubechies.maskPoly, Daubechies.coeff, Daubechies.coeff_natCast, Daubechies.symbol}
\leanok
\uses{def:daub-mask}
Let $L$ be a Daubechies mask of order $N$. Its \emph{mask polynomial} is
\[
  m(z) = \Bigl(\frac{1 + z}{2}\Bigr)^N L(z) .
\]
Its \emph{filter coefficients} are $c_k = 2 [z^k] m$ for $k \in \N$, where $[z^k] m$ is the
coefficient of $z^k$ in $m$, and $c_k = 0$ for the integers $k < 0$. Its \emph{symbol} is the
function $m_0(\xi) = m(e^{-2\pi i \xi})$ of $\xi \in \R$.
\end{definition}

The filter coefficients, the symbol and everything built from them below belong to a Daubechies
mask $L$ of order $N$, which the notation omits.

\begin{lemma}[The mask polynomial]
\label{lem:daub-mask-poly}
\lean{Daubechies.X_add_one_pow_dvd_maskPoly, Daubechies.natDegree_maskPoly_le,
  Daubechies.natDegree_maskPoly_lt, Daubechies.eval_one_maskPoly, Daubechies.coeff_eq_zero,
  Daubechies.sum_coeff}
\leanok
\uses{def:daub-symbol}
For every Daubechies mask of order $N$ the mask polynomial $m$ is divisible by $(z + 1)^N$, has
degree at most $2N - 1$, and satisfies $m(1) = 1$. The filter coefficients satisfy $c_k = 0$ for
every integer $k < 0$ and every integer $k \geq 2N$, and
\[
  \sum_{k=0}^{2N-1} c_k = 2 .
\]
\end{lemma}

\begin{proof}
\leanok
$m(z) = (z + 1)^N 2^{-N} L(z)$ is the product of $N$ factors of degree $1$ and of $L$, of degree
at most $N - 1$, and $m(1) = L(1) = 1$. The coefficients of $m$ add up to $m(1) = 1$, and
$c_k = 2 [z^k] m$.
\end{proof}

\begin{lemma}[The symbol as a trigonometric polynomial]
\label{lem:daub-symbol-sum}
\lean{Daubechies.symbol_eq_mul, Daubechies.symbol_eq_sum, Daubechies.symbol_zero}
\leanok
\uses{def:daub-symbol}
For every Daubechies mask of order $N$ and every $\xi \in \R$,
\[
  m_0(\xi) = \Bigl(\frac{1 + e^{-2\pi i \xi}}{2}\Bigr)^N L(e^{-2\pi i \xi})
  = \frac12 \sum_{k=0}^{2N-1} c_k e^{-2\pi i k \xi} ,
\]
and $m_0(0) = 1$.
\end{lemma}

\begin{proof}
\leanok
\uses{lem:daub-mask-poly, lem:daub-exp}
By \Cref{lem:daub-mask-poly}, $m(z) = \frac12 \sum_{k < 2N} c_k z^k$. Evaluate at
$z = e^{-2\pi i \xi}$, where $z^k = e^{-2\pi i k \xi}$ (\Cref{lem:daub-exp}). At $\xi = 0$,
$z = 1$ and $m(1) = 1$.
\end{proof}

\begin{proposition}[The modulus of the symbol]
\label{prop:daub-symbol-modulus}
\lean{Daubechies.normSq_symbol}
\leanok
\uses{def:daub-symbol}
For every Daubechies mask of order $N$ and every $\xi \in \R$,
\[
  |m_0(\xi)|^2 = \cos(\pi\xi)^{2N} P(\xi) .
\]
\end{proposition}

\begin{proof}
\leanok
\uses{lem:daub-symbol-sum, lem:daub-half-factor, def:daub-mask}
By \Cref{lem:daub-symbol-sum}, $|m_0(\xi)|^2$ is the product of
$\bigl(|1 + e^{-2\pi i \xi}|^2/4\bigr)^N = \cos(\pi\xi)^{2N}$ (\Cref{lem:daub-half-factor}) and
$|L(e^{-2\pi i \xi})|^2 = P(\xi)$ (\Cref{def:daub-mask}).
\end{proof}

\begin{theorem}[The quadrature mirror filter identity]
\label{thm:daub-qmf}
\lean{Daubechies.qmf, Daubechies.norm_symbol_le_one}
\leanok
\uses{def:daub-symbol}
For every Daubechies mask and every $\xi \in \R$,
\[
  |m_0(\xi)|^2 + \bigl|m_0\bigl(\xi + \tfrac12\bigr)\bigr|^2 = 1 ;
\]
in particular $|m_0(\xi)| \leq 1$.
\end{theorem}

\begin{proof}
\leanok
\uses{prop:daub-symbol-modulus, thm:daub-fun-bezout}
By \Cref{prop:daub-symbol-modulus} at $\xi$ and at $\xi + \frac12$, where
$\cos(\pi\xi + \frac{\pi}{2})^{2N} = \sin(\pi\xi)^{2N}$, the left side is
$\cos(\pi\xi)^{2N} P(\xi) + \sin(\pi\xi)^{2N} P(\xi + \frac12)$. This is $1$ by the
trigonometric B\'ezout identity (\Cref{thm:daub-fun-bezout}), since $N \geq 1$.
\end{proof}

\begin{lemma}[Symmetries and zeros of the symbol]
\label{lem:daub-symbol-props}
\lean{Daubechies.symbol_periodic, Daubechies.symbol_neg, Daubechies.symbol_ne_zero,
  Daubechies.symbol_half_add_int, Daubechies.symbol_half, Daubechies.contDiff_symbol,
  Daubechies.continuous_symbol}
\leanok
\uses{def:daub-symbol}
For every Daubechies mask the symbol $m_0$ is smooth and has period $1$, and
$m_0(-\xi) = \overline{m_0(\xi)}$ for every $\xi \in \R$. Moreover $m_0(\xi) \neq 0$ for every
$\xi \in \R$ with $|\xi| < \frac12$, and $m_0(\frac12 + n) = 0$ for every $n \in \Z$.
\end{lemma}

\begin{proof}
\leanok
\uses{lem:daub-exp, lem:daub-symbol-sum, prop:daub-symbol-modulus, prop:daub-fun-bounds}
The symbol is a polynomial with real coefficients in $e^{-2\pi i \xi}$, which is smooth, has
period $1$ and satisfies $\overline{e^{-2\pi i \xi}} = e^{-2\pi i (-\xi)}$ (\Cref{lem:daub-exp}).
For $|\xi| < \frac12$, $\cos \pi\xi > 0$ and $P(\xi) \geq 1$ (\Cref{prop:daub-fun-bounds}), so
$|m_0(\xi)|^2 > 0$ by \Cref{prop:daub-symbol-modulus}. At $\xi = \frac12 + n$,
$z = e^{-2\pi i \xi} = -1$, where the factor $((1 + z)/2)^N$ of \Cref{lem:daub-symbol-sum}
vanishes since $N \geq 1$.
\end{proof}

\begin{lemma}[The symbol near $0$]
\label{lem:daub-symbol-lipschitz}
\lean{Daubechies.symbolLip, Daubechies.symbolLip_nonneg, Daubechies.norm_symbol_sub_one_le}
\leanok
\uses{def:daub-symbol}
For a Daubechies mask of order $N$ let $C_0 = \pi \sum_{k=0}^{2N-1} k |c_k|$. Then $C_0 \geq 0$,
and for every $\xi \in \R$
\[
  |m_0(\xi) - 1| \leq C_0 |\xi| .
\]
\end{lemma}

\begin{proof}
\leanok
\uses{lem:daub-mask-poly, lem:daub-symbol-sum, lem:daub-exp}
Since $\sum_k c_k = 2$ (\Cref{lem:daub-mask-poly}), \Cref{lem:daub-symbol-sum} gives
$m_0(\xi) - 1 = \frac12 \sum_k c_k (e^{-2\pi i k \xi} - 1)$, and
$|e^{-2\pi i k \xi} - 1| \leq 2\pi k |\xi|$ (\Cref{lem:daub-exp}).
\end{proof}

\paragraph{Sums over $\Z$.}
The identities below are sums over $\Z$ with finitely many nonzero terms (Mathlib's
\texttt{finsum}). Throughout, $\delta_{ab} = 1$ if $a = b$ and $\delta_{ab} = 0$ otherwise.

\begin{lemma}[Sums over $\Z$ by parity]
\label{lem:daub-parity-split}
\lean{Daubechies.finsum_eq_sum_range_of_eq_zero,
  Daubechies.finsum_two_mul_add_finsum_two_mul_add_one, Daubechies.finsum_two_mul_add_add_finsum,
  Daubechies.finsum_two_mul_add_add_two_mul, Daubechies.neg_one_zpow_eq_of_even_sub,
  Daubechies.neg_one_zpow_eq_neg_of_odd_sub, Daubechies.finsum_two_mul_add_eq_half}
\leanok
Let $g$ be a function on $\Z$ with values in a commutative monoid, written additively, and let
$n \in \N$ and $r, q \in \Z$.
\begin{enumerate}
  \item If $g(k) = 0$ for every $k < 0$ and every $k \geq n$, then
    $\sum_{k \in \Z} g(k) = \sum_{k=0}^{n-1} g(k)$.
  \item $\sum_{i \in \Z} g(2i + r + 2q) = \sum_{i \in \Z} g(2i + r)$, and if $g$ has finite
    support, then
    \[
      \sum_{i \in \Z} g(2i + r) + \sum_{i \in \Z} g(2i + r + 1) = \sum_{k \in \Z} g(k) .
    \]
  \item If $g$ is real valued with finite support and $\sum_{k \in \Z} (-1)^k g(k) = 0$, then
    $\sum_{i \in \Z} g(2i + r) = \frac12 \sum_{k \in \Z} g(k)$.
\end{enumerate}
For all integers $a$ and $b$, $(-1)^a = (-1)^b$ if $a - b$ is even, and $(-1)^a = -(-1)^b$ if
$a - b$ is odd.
\end{lemma}

\begin{proof}
\leanok
The even and the odd integers partition $\Z$, and $k \mapsto k + r$ and $i \mapsto i + q$ are
bijections of $\Z$. In item~3 the alternating sum is
$(-1)^r \bigl(\sum_i g(2i + r) - \sum_i g(2i + r + 1)\bigr)$, so the two parity classes carry the
same sum, and by item~2 the two sums add up to $\sum_k g(k)$.
\end{proof}

\begin{lemma}[The filter coefficients as a sequence on $\Z$]
\label{lem:daub-coeff-finsum}
\lean{Daubechies.finite_support_coeff, Daubechies.finsum_eq_sum_range_of_coeff,
  Daubechies.finsum_coeff, Daubechies.symbol_eq_finsum}
\leanok
\uses{def:daub-symbol}
For every Daubechies mask of order $N$ the sequence $(c_k)_{k \in \Z}$ has finite support,
$\sum_{k \in \Z} f(k) = \sum_{k=0}^{2N-1} f(k)$ for every function $f$ on $\Z$ that vanishes
wherever $c_k = 0$, and $\sum_{k \in \Z} c_k = 2$. For every $\xi \in \R$,
\[
  m_0(\xi) = \frac12 \sum_{k \in \Z} c_k e^{-2\pi i k \xi} .
\]
\end{lemma}

\begin{proof}
\leanok
\uses{lem:daub-mask-poly, lem:daub-parity-split, lem:daub-symbol-sum}
$c_k = 0$ for $k < 0$ and for $k \geq 2N$ (\Cref{lem:daub-mask-poly}); apply item~1 of
\Cref{lem:daub-parity-split}, and \Cref{lem:daub-symbol-sum}.
\end{proof}

\paragraph{Orthogonality of the even shifts.}
The coefficients of a trigonometric polynomial are its Fourier coefficients on $[0, 1]$, so
integrating the quadrature mirror filter identity against exponentials gives identities for the
filter coefficients.

\begin{lemma}[Orthonormality of the exponentials]
\label{lem:daub-exp-orthonormal}
\lean{Daubechies.continuous_expFreq_mul, Daubechies.integral_expFreq_intCast_mul,
  Daubechies.integral_sum_mul_conj_sum, Daubechies.integral_finsum_mul_conj_finsum,
  Daubechies.integral_expFreq_mul_normSq_finsum}
\leanok
For every $c \in \R$ the function $\xi \mapsto e^{-2\pi i c \xi}$ is continuous, and for every
$n \in \Z$, $\int_0^1 e^{-2\pi i n \xi} \, d\xi = \delta_{n0}$. For every finite set
$u \subseteq \Z$ and every $a : \Z \to \C$ let
$T_a(\xi) = \sum_{k \in u} a_k e^{-2\pi i k \xi}$. Then for all $a, b : \Z \to \C$
\[
  \int_0^1 T_a(\xi) \, \overline{T_b(\xi)} \, d\xi = \sum_{k \in u} a_k \overline{b_k} ,
\]
and the same holds for the sums over $k \in \Z$ if $a$ and $b$ have finite support. For every
$a : \Z \to \C$ with finite support and every $n \in \Z$,
\[
  \int_0^1 e^{-2\pi i n \xi} \Bigl|\sum_{k \in \Z} a_k e^{-2\pi i k \xi}\Bigr|^2 \, d\xi
  = \sum_{k \in \Z} a_{k-n} \overline{a_k} .
\]
\end{lemma}

\begin{proof}
\leanok
\uses{lem:daub-exp}
For $n \neq 0$, $\int_0^1 e^{-2\pi i n \xi} \, d\xi = (e^{-2\pi i n} - 1)/(-2\pi i n) = 0$.
Expanding the product gives the double sum of the terms
$a_k \overline{b_l} e^{-2\pi i (k - l) \xi}$ (\Cref{lem:daub-exp}), of which only those with
$k = l$ have nonzero integrals. For the last identity, multiplication by $e^{-2\pi i n \xi}$
shifts the coefficients to $a_{k-n}$, and $|z|^2 = z \bar z$.
\end{proof}

\begin{lemma}[Fourier coefficients of $|m_0|^2$]
\label{lem:daub-symbol-fourier}
\lean{Daubechies.norm_half_mul_sq, Daubechies.integral_expFreq_mul_normSq_symbol,
  Daubechies.integral_expFreq_mul_normSq_symbol_add_half}
\leanok
\uses{def:daub-symbol}
For every Daubechies mask and every $n \in \Z$,
\[
  \int_0^1 e^{-2\pi i n \xi} |m_0(\xi)|^2 \, d\xi = \frac14 \sum_{k \in \Z} c_{k-n} c_k
\]
and
\[
  \int_0^1 e^{-2\pi i n \xi} \bigl|m_0\bigl(\xi + \tfrac12\bigr)\bigr|^2 \, d\xi
  = \frac14 e^{\pi i n} \sum_{k \in \Z} c_{k-n} c_k .
\]
\end{lemma}

\begin{proof}
\leanok
\uses{lem:daub-coeff-finsum, lem:daub-exp-orthonormal}
By \Cref{lem:daub-coeff-finsum}, $m_0(\xi) = \frac12 \sum_k c_k e^{-2\pi i k \xi}$ and
$m_0(\xi + \frac12) = \frac12 \sum_k c_k e^{-\pi i k} e^{-2\pi i k \xi}$; apply the last identity
of \Cref{lem:daub-exp-orthonormal}, with $|z/2|^2 = |z|^2/4$. In the second case the factors
combine to $e^{-\pi i (k - n)} \, \overline{e^{-\pi i k}} = e^{\pi i n}$.
\end{proof}

\begin{theorem}[Orthogonality of the even shifts]
\label{thm:daub-even-shifts}
\lean{Daubechies.sum_coeff_mul_coeff_add}
\leanok
\uses{def:daub-symbol}
For every Daubechies mask and every $j \in \Z$,
\[
  \sum_{k \in \Z} c_k c_{k+2j} = 2 \delta_{j0} .
\]
\end{theorem}

\begin{proof}
\leanok
\uses{thm:daub-qmf, lem:daub-symbol-fourier, lem:daub-exp-orthonormal}
Multiply the quadrature mirror filter identity (\Cref{thm:daub-qmf}) by $e^{2\pi i \cdot 2j \xi}$
and integrate over $[0, 1]$. By \Cref{lem:daub-symbol-fourier} with $n = -2j$, where
$e^{\pi i n} = 1$, the left side becomes $\frac12 \sum_k c_k c_{k+2j}$, and the right side is
$\delta_{j0}$ by \Cref{lem:daub-exp-orthonormal}.
\end{proof}

\paragraph{Sum rules.}
The factor $(z + 1)^N$ of the mask polynomial makes the alternating moments of order less than
$N$ of the filter coefficients vanish. For a real polynomial $p$ let $E p = z p'$ be the Euler
operator.

\begin{lemma}[The Euler operator]
\label{lem:daub-euler-operator}
\lean{Daubechies.coeff_X_mul_derivative, Daubechies.coeff_iterate_X_mul_derivative,
  Daubechies.pow_dvd_iterate_X_mul_derivative,
  Daubechies.sum_range_neg_one_pow_mul_pow_mul_coeff}
\leanok
Let $p$ be a real polynomial and $j, k, l, n \in \N$. The coefficient of $z^k$ in $E^l p$ is
$k^l$ times the coefficient of $z^k$ in $p$. If $(z + 1)^j$ divides $p$ and $l \leq j$, then
$(z + 1)^{j-l}$ divides $E^l p$. If $(z + 1)^j$ divides $p$, $l < j$ and $p$ has degree less
than $n$, then
\[
  \sum_{k=0}^{n-1} (-1)^k k^l [z^k] p = 0 .
\]
\end{lemma}

\begin{proof}
\leanok
$E z^k = k z^k$. Differentiation lowers the power of $z + 1$ dividing a polynomial by at most
one, and multiplication by $z$ keeps it. The sum is the value of $E^l p$ at $z = -1$, a zero of
$E^l p$ since $j - l \geq 1$.
\end{proof}

\begin{theorem}[The sum rules]
\label{thm:daub-sum-rules}
\lean{Daubechies.sum_range_rule, Daubechies.sum_rule, Daubechies.sum_rule_poly}
\leanok
\uses{def:daub-symbol}
For every Daubechies mask of order $N$ and every real polynomial $Q$ of degree less than $N$,
\[
  \sum_{k \in \Z} (-1)^k Q(k) \, c_k = 0 .
\]
In particular $\sum_{k \in \Z} (-1)^k k^l c_k = \sum_{k=0}^{2N-1} (-1)^k k^l c_k = 0$ for every
$l < N$.
\end{theorem}

\begin{proof}
\leanok
\uses{lem:daub-euler-operator, lem:daub-mask-poly, lem:daub-coeff-finsum}
\Cref{lem:daub-euler-operator} for $p = m$, $j = N$ and $n = 2N$ (\Cref{lem:daub-mask-poly}),
with $c_k = 2 [z^k] m$, gives the case $Q(k) = k^l$. The sum over $\Z$ is the sum over
$0 \leq k < 2N$ (\Cref{lem:daub-coeff-finsum}), and $Q$ is a linear combination of the monomials
of degree less than $N$.
\end{proof}

\begin{corollary}[Parity sums]
\label{cor:daub-parity-sums}
\lean{Daubechies.finsum_neg_one_zpow_mul_coeff, Daubechies.sum_coeff_two_mul_add,
  Daubechies.sum_coeff_two_mul, Daubechies.sum_coeff_two_mul_add_one}
\leanok
\uses{def:daub-symbol}
For every Daubechies mask and every integer $r$,
\[
  \sum_{k \in \Z} (-1)^k c_k = 0 \quad\text{and}\quad \sum_{i \in \Z} c_{2i+r} = 1 ;
\]
in particular $\sum_{i \in \Z} c_{2i} = \sum_{i \in \Z} c_{2i+1} = 1$.
\end{corollary}

\begin{proof}
\leanok
\uses{thm:daub-sum-rules, lem:daub-parity-split, lem:daub-coeff-finsum}
The first identity is \Cref{thm:daub-sum-rules} for $l = 0 < N$. Item~3 of
\Cref{lem:daub-parity-split} and $\sum_k c_k = 2$ (\Cref{lem:daub-coeff-finsum}) give the second.
\end{proof}

The first identity says $m_0(\frac12) = 0$.

\begin{proposition}[The first moment]
\label{prop:daub-first-moments}
\lean{Daubechies.firstMoment, Daubechies.sum_two_mul_add_mul_coeff}
\leanok
\uses{def:daub-symbol}
For a Daubechies mask let $\mu = \frac12 \sum_{k \in \Z} k c_k$ be its \emph{first moment}. If
the mask has order $N \geq 2$, then for every $r \in \Z$
\[
  \sum_{i \in \Z} (2i + r) \, c_{2i+r} = \mu .
\]
\end{proposition}

\begin{proof}
\leanok
\uses{thm:daub-sum-rules, lem:daub-parity-split, lem:daub-coeff-finsum}
Item~3 of \Cref{lem:daub-parity-split} for $g(k) = k c_k$, which has finite support
(\Cref{lem:daub-coeff-finsum}) and whose alternating sum vanishes by \Cref{thm:daub-sum-rules}
with $l = 1 < N$.
\end{proof}

\begin{lemma}[Derivatives of the symbol]
\label{lem:daub-symbol-derivatives}
\lean{Daubechies.iteratedDeriv_sum_mul_cexp, Daubechies.symbol_eq_sum_cexp,
  Daubechies.iteratedDeriv_symbol}
\leanok
\uses{def:daub-symbol}
For every finite family of pairs $a_i, b_i \in \C$ and every $l \in \N$, the $l$-th derivative
of $t \mapsto \sum_i a_i e^{b_i t}$ on $\R$ is $t \mapsto \sum_i a_i b_i^l e^{b_i t}$. For every
Daubechies mask of order $N$, every $l \in \N$ and every $\xi \in \R$,
\[
  m_0^{(l)}(\xi) = \frac12 \sum_{k=0}^{2N-1} c_k (-2\pi i k)^l e^{-2\pi i k \xi} .
\]
\end{lemma}

\begin{proof}
\leanok
\uses{lem:daub-symbol-sum}
Induction on $l$, with $\frac{d}{dt} e^{bt} = b e^{bt}$, and
$m_0(\xi) = \frac12 \sum_k c_k e^{-2\pi i k \xi}$ (\Cref{lem:daub-symbol-sum}).
\end{proof}

\begin{proposition}[Zeros of the symbol at the half-integers]
\label{prop:daub-symbol-zeros}
\lean{Daubechies.iteratedDeriv_symbol_half_add_int}
\leanok
\uses{def:daub-symbol}
For every Daubechies mask of order $N$, every $l < N$ and every $n \in \Z$,
$m_0^{(l)}(\frac12 + n) = 0$.
\end{proposition}

\begin{proof}
\leanok
\uses{lem:daub-symbol-derivatives, lem:daub-exp, thm:daub-sum-rules}
At $\xi = \frac12 + n$, $e^{-2\pi i k \xi} = (-1)^k$ (\Cref{lem:daub-exp}), so
$m_0^{(l)}(\frac12 + n) = \frac12 (-2\pi i)^l \sum_k (-1)^k k^l c_k$ by
\Cref{lem:daub-symbol-derivatives}, which vanishes by \Cref{thm:daub-sum-rules}.
\end{proof}

\paragraph{The dual filter.}
The dual filter carries the coefficients of the wavelet. Its identities follow from those of
the filter by reflections $k \mapsto a - k$ of the index and a split by parity.

\begin{definition}[The dual filter]
\label{def:daub-dual-filter}
\lean{Daubechies.dualCoeff, Daubechies.dualCoeff_eq_zero, Daubechies.finite_support_dualCoeff,
  Daubechies.finsum_eq_sum_Icc_of_dualCoeff}
\leanok
\uses{def:daub-symbol}
The \emph{dual filter} of a Daubechies mask of order $N$ is $g_k = (-1)^k c_{1-k}$, $k \in \Z$.
It vanishes outside $2 - 2N \leq k \leq 1$, so it has finite support, and
$\sum_{k \in \Z} f(k) = \sum_{k=2-2N}^{1} f(k)$ for every function $f$ on $\Z$ that vanishes
wherever $g_k = 0$.
\end{definition}

\begin{proposition}[Orthogonality of the dual filter]
\label{prop:daub-dual-orthogonality}
\lean{Daubechies.sum_dualCoeff_mul_dualCoeff_add, Daubechies.sum_coeff_mul_dualCoeff_add}
\leanok
\uses{def:daub-dual-filter}
For every Daubechies mask and every $j \in \Z$,
\[
  \sum_{k \in \Z} g_k g_{k+2j} = 2 \delta_{j0} \quad\text{and}\quad
  \sum_{k \in \Z} c_k g_{k+2j} = 0 .
\]
\end{proposition}

\begin{proof}
\leanok
\uses{thm:daub-even-shifts, lem:daub-parity-split}
The substitution $k \mapsto 1 - 2j - k$ turns the first sum into $\sum_k c_k c_{k+2j}$, since
$1 - 2j - k$ and $1 - k$ have the same parity (\Cref{lem:daub-parity-split}); apply
\Cref{thm:daub-even-shifts}. It turns the second sum into its negative, since $1 - k$ and
$k + 2j$ have different parity.
\end{proof}

\begin{lemma}[Reindexing over one parity class]
\label{lem:daub-dual-reindex}
\lean{Daubechies.finsum_coeff_sub_two_mul_mul, Daubechies.finsum_dualCoeff_sub_two_mul_mul}
\leanok
\uses{def:daub-dual-filter}
For every Daubechies mask and all $a, b \in \Z$,
\[
  \sum_{j \in \Z} c_{a-2j} c_{b-2j} = \sum_{i \in \Z} c_{2i+b+(a-b)} c_{2i+b}
\]
and
\[
  \sum_{j \in \Z} g_{a-2j} g_{b-2j}
  = (-1)^{a+b} \sum_{i \in \Z} c_{2i+(1-a)+(a-b)} c_{2i+(1-a)} .
\]
\end{lemma}

\begin{proof}
\leanok
\uses{lem:daub-parity-split}
In the first sum substitute $j = -i$. In the second,
$(-1)^{a-2j} (-1)^{b-2j} = (-1)^{a+b}$ (\Cref{lem:daub-parity-split}),
$1 - (a - 2j) = 2j + (1 - a)$ and $1 - (b - 2j) = 2j + (1 - a) + (a - b)$.
\end{proof}

\begin{theorem}[The filter and its dual resolve the identity]
\label{thm:daub-filter-identity}
\lean{Daubechies.sum_coeff_mul_add_dual}
\leanok
\uses{def:daub-dual-filter}
For every Daubechies mask and all $a, b \in \Z$,
\[
  \sum_{j \in \Z} \bigl(c_{a-2j} c_{b-2j} + g_{a-2j} g_{b-2j}\bigr) = 2 \delta_{ab} .
\]
\end{theorem}

\begin{proof}
\leanok
\uses{lem:daub-dual-reindex, lem:daub-parity-split, thm:daub-even-shifts}
By \Cref{lem:daub-dual-reindex} both sums run over one parity class of
$h(k) = c_{k+(a-b)} c_k$, the first over the class of $b$ and the second over the class of
$1 - a$. If $a - b$ is even, then $(-1)^{a+b} = 1$ and the two classes are complementary, so the
sum is $\sum_k c_k c_{k+(a-b)} = 2 \delta_{ab}$ (\Cref{lem:daub-parity-split},
\Cref{thm:daub-even-shifts}). If $a - b$ is odd, the two classes agree and $(-1)^{a+b} = -1$, so
the two sums cancel.
\end{proof}

\begin{proposition}[Vanishing moments of the dual filter]
\label{prop:daub-dual-moments}
\lean{Daubechies.sum_dualCoeff_mul_pow}
\leanok
\uses{def:daub-dual-filter}
For every Daubechies mask of order $N$ and every $l < N$, $\sum_{k \in \Z} g_k k^l = 0$.
\end{proposition}

\begin{proof}
\leanok
\uses{thm:daub-sum-rules}
The substitution $k \mapsto 1 - k$ turns the sum into $-\sum_k (-1)^k (1 - k)^l c_k$, which
vanishes by \Cref{thm:daub-sum-rules} for $Q(z) = (1 - z)^l$.
\end{proof}

\paragraph{The mask D4.}
For $N = 2$ the polynomial $L$ is linear, with its root $2 + \sqrt 3$ outside the unit disk.

\begin{lemma}[Linear polynomials on the unit circle]
\label{lem:daub-linear-circle}
\lean{Daubechies.normSq_add_mul_expFreq, Daubechies.daubFun_two_eq_two_sub_cos}
\leanok
\uses{def:daub-fun}
For all $a, b, \xi \in \R$,
\[
  |a + b e^{-2\pi i \xi}|^2 = a^2 + b^2 + 2ab \cos 2\pi\xi .
\]
For $N = 2$ and every $\xi \in \R$, $P(\xi) = 2 - \cos 2\pi\xi$.
\end{lemma}

\begin{proof}
\leanok
\uses{lem:daub-exp}
$|a + bz|^2 = a^2 + b^2 |z|^2 + 2ab \operatorname{Re} z$ for real $a$ and $b$, with $|z| = 1$ and
$\operatorname{Re} z = \cos 2\pi\xi$ at $z = e^{-2\pi i \xi}$ (\Cref{lem:daub-exp}). For $N = 2$,
$P(\xi) = 1 + 2 \sin^2 \pi\xi$ (\Cref{def:daub-fun}) and $\cos 2\pi\xi = 1 - 2 \sin^2 \pi\xi$.
\end{proof}

\begin{proposition}[The mask D4]
\label{prop:daub-d4}
\lean{Daubechies.d4, Daubechies.d4_L}
\leanok
\uses{def:daub-mask}
The polynomial
\[
  L(z) = \frac{1 + \sqrt 3}{2} + \frac{1 - \sqrt 3}{2} \, z
\]
is a Daubechies mask of order $2$, the mask \emph{D4}.
\end{proposition}

\begin{proof}
\leanok
\uses{lem:daub-linear-circle}
$L$ has degree $1 < 2$ and $L(1) = 1$. With $a = (1 + \sqrt 3)/2$ and $b = (1 - \sqrt 3)/2$,
$a^2 + b^2 = 2$ and $2ab = -1$, so $|L(e^{-2\pi i \xi})|^2 = 2 - \cos 2\pi\xi = P(\xi)$ by
\Cref{lem:daub-linear-circle}.
\end{proof}

\begin{proposition}[The filter coefficients of D4]
\label{prop:daub-d4-coeff}
\lean{Daubechies.maskPoly_d4, Daubechies.coeff_d4_zero, Daubechies.coeff_d4_one,
  Daubechies.coeff_d4_two, Daubechies.coeff_d4_three, Daubechies.coeff_d4,
  Daubechies.coeff_d4_eq_zero}
\leanok
\uses{prop:daub-d4, def:daub-symbol}
The mask polynomial of D4 is
\[
  m(z) = \frac18 \bigl((1 + \sqrt 3) + (3 + \sqrt 3) z + (3 - \sqrt 3) z^2
    + (1 - \sqrt 3) z^3\bigr) ,
\]
its filter coefficients are
\[
  c_0 = \frac{1 + \sqrt 3}{4} , \quad c_1 = \frac{3 + \sqrt 3}{4} , \quad
  c_2 = \frac{3 - \sqrt 3}{4} , \quad c_3 = \frac{1 - \sqrt 3}{4} ,
\]
and $c_k = 0$ for every other $k \in \Z$.
\end{proposition}

\begin{proof}
\leanok
\uses{lem:daub-mask-poly}
Expand $m(z) = \frac14 (1 + 2z + z^2) L(z)$; two real polynomials agree once they agree at every
real point. Then $c_k = 2 [z^k] m$, and $c_k = 0$ for $k < 0$ and for $k \geq 4$
(\Cref{lem:daub-mask-poly}).
\end{proof}

The orthonormal filter $c_k/\sqrt 2$ is the form in which these coefficients are usually
tabulated.

\section{The symbol product}

The scaling function of a mask is the inverse Fourier transform of the infinite product
$\Phi(\xi) = \prod_{j \geq 0} m_0(\xi/2^{j+1})$. This section proves that the product converges
locally uniformly, satisfies $\Phi(2\xi) = m_0(\xi) \Phi(\xi)$, and has the modulus
$|\Phi(\xi)|^2 = \operatorname{sinc}(\pi\xi)^{2N} \prod_{j \geq 0} P(\xi/2^{j+1})$.

\paragraph{The cosine product.}
Infinite products are understood as in Mathlib: $\prod_{j \geq 0} a_j$ converges if the products
$\prod_{j \in F} a_j$ over the finite sets $F \subseteq \N$ converge as $F$ increases, and then
the partial products $\prod_{j<n} a_j$ converge to the same value. Here
$\operatorname{sinc} x = \sin x / x$ for $x \neq 0$ and $\operatorname{sinc} 0 = 1$. The doubling
formula $\sin 2y = 2 \sin y \cos y$ telescopes into the cosine product, which for $x = \pi/2$ is
Vi\`ete's product $2/\pi = \prod_{j \geq 0} \cos(\pi/2^{j+2})$.

\begin{lemma}[Dyadic scalings]
\label{lem:daub-dyadic}
\lean{Daubechies.div_two_pow_add_succ, Daubechies.sq_div_two_pow_succ,
  Daubechies.hasSum_sq_div_two_pow}
\leanok
For every $x \in \R$ and all $n, j \in \N$, $x/2^{n+j+1} = (x/2^n)/2^{j+1}$ and
$(x/2^{j+1})^2 = \frac{x^2}{4} \, 4^{-j}$, and
\[
  \sum_{j \geq 0} \Bigl(\frac{x}{2^{j+1}}\Bigr)^2 = \frac{x^2}{3} .
\]
\end{lemma}

\begin{proof}
\leanok
A geometric series: $\frac{x^2}{4} \sum_{j \geq 0} 4^{-j} = \frac{x^2}{4} \cdot \frac43$.
\end{proof}

\begin{lemma}[The finite cosine product]
\label{lem:daub-finite-cos-product}
\lean{Daubechies.sinc_mul_cos, Daubechies.sinc_div_two_pow_mul_prod_cos,
  Daubechies.norm_cos_sub_one_le, Daubechies.multipliable_cos_div_two_pow}
\leanok
For every $x \in \R$ and every $n \in \N$,
\[
  \operatorname{sinc}\Bigl(\frac{x}{2^n}\Bigr) \prod_{j<n} \cos\Bigl(\frac{x}{2^{j+1}}\Bigr)
  = \operatorname{sinc} x ,
\]
and the product $\prod_{j \geq 0} \cos(x/2^{j+1})$ converges. For every $y \in \R$,
$\operatorname{sinc}(y) \cos(y) = \operatorname{sinc}(2y)$ and $|\cos y - 1| \leq y^2/2$.
\end{lemma}

\begin{proof}
\leanok
\uses{lem:daub-dyadic}
$\sin 2y = 2 \sin y \cos y$, and $1 - y^2/2 \leq \cos y \leq 1$. Induction on $n$ with
$y = x/2^{n+1}$ gives the identity. The factors are $1 + a_j$ with
$|a_j| \leq (x/2^{j+1})^2/2$, a summable bound (\Cref{lem:daub-dyadic}).
\end{proof}

\begin{theorem}[The cosine product]
\label{thm:daub-cos-product}
\lean{Daubechies.hasProd_cos_div_two_pow, Daubechies.tprod_cos}
\leanok
For every $x \in \R$,
\[
  \prod_{j \geq 0} \cos\Bigl(\frac{x}{2^{j+1}}\Bigr) = \operatorname{sinc} x .
\]
\end{theorem}

\begin{proof}
\leanok
\uses{lem:daub-finite-cos-product}
The partial products converge (\Cref{lem:daub-finite-cos-product}), and
$\operatorname{sinc}(x/2^n) \to \operatorname{sinc} 0 = 1$, so by the identity of
\Cref{lem:daub-finite-cos-product} their limit is $\operatorname{sinc} x$.
\end{proof}

\paragraph{The product of $P$.}
The modulus of $\Phi$ contains the real product $\prod_{j \geq 0} P(\xi/2^{j+1})$; its tail after
the first $n$ factors is bounded uniformly for $|\xi| \leq 2^n/2$.

\begin{lemma}[Convergence of the product of $P$]
\label{lem:daub-fun-product}
\lean{Daubechies.multipliable_daubFun, Daubechies.multipliable_daubFun_shift,
  Daubechies.tprod_daubFun_nonneg, Daubechies.tprod_daubFun_eq_prod_mul}
\leanok
\uses{def:daub-fun}
For every $N \in \N$, every $n \in \N$ and every $\xi \in \R$ the products
$\prod_{j \geq 0} P(\xi/2^{j+1})$ and $\prod_{j \geq 0} P(\xi/2^{n+j+1})$ converge, the first
one is nonnegative, and
\[
  \prod_{j \geq 0} P\Bigl(\frac{\xi}{2^{j+1}}\Bigr)
  = \prod_{j<n} P\Bigl(\frac{\xi}{2^{j+1}}\Bigr) \cdot
    \prod_{j \geq 0} P\Bigl(\frac{\xi}{2^{n+j+1}}\Bigr) .
\]
\end{lemma}

\begin{proof}
\leanok
\uses{prop:daub-fun-bounds, lem:daub-dyadic}
For $N \geq 1$ the factors are $1 + a_j$ with
$0 \leq a_j \leq \binom{2N-1}{N} \pi^2 (\xi/2^{j+1})^2$ (\Cref{prop:daub-fun-bounds}), a
summable bound (\Cref{lem:daub-dyadic}); for $N = 0$ every factor is $0$. The second product is
the first one at the point $\xi/2^n$ (\Cref{lem:daub-dyadic}), and $P \geq 0$.
\end{proof}

\begin{proposition}[Bounds for the product of $P$]
\label{prop:daub-fun-product-bound}
\lean{Daubechies.tprod_daubFun_le, Daubechies.tprod_daubFun_shift_le}
\leanok
\uses{def:daub-fun}
For every $N \in \N$ and every $x \in \R$,
\[
  \prod_{j \geq 0} P\Bigl(\frac{x}{2^{j+1}}\Bigr)
  \leq \exp\Bigl(\binom{2N-1}{N} \frac{\pi^2 x^2}{3}\Bigr) .
\]
For every $n \in \N$ and every $\xi \in \R$ with $|\xi| \leq 2^n/2$,
\[
  \prod_{j \geq 0} P\Bigl(\frac{\xi}{2^{n+j+1}}\Bigr)
  \leq \exp\Bigl(\binom{2N-1}{N} \frac{\pi^2}{12}\Bigr) .
\]
\end{proposition}

\begin{proof}
\leanok
\uses{prop:daub-fun-bounds, lem:daub-dyadic, lem:daub-fun-product}
With $t_j = \binom{2N-1}{N} \pi^2 (x/2^{j+1})^2$, each factor is at most $1 + t_j \leq e^{t_j}$
(\Cref{prop:daub-fun-bounds}), and $\sum_j (x/2^{j+1})^2 = x^2/3$ (\Cref{lem:daub-dyadic}). The
tail is the product at $x = \xi/2^n$ (\Cref{lem:daub-fun-product}), where $|x| \leq \frac12$
gives $x^2/3 \leq 1/12$.
\end{proof}

\paragraph{Convergence.}
Near $0$ the symbol is $1 + O(|\eta|)$, so the factors of $\Phi$ approach $1$ geometrically.

\begin{definition}[The symbol product]
\label{def:daub-symbol-product}
\lean{Daubechies.symbolProduct}
\leanok
\uses{def:daub-symbol}
The \emph{symbol product} of a Daubechies mask is the function
\[
  \Phi(\xi) = \prod_{j \geq 0} m_0\Bigl(\frac{\xi}{2^{j+1}}\Bigr) , \qquad \xi \in \R .
\]
\end{definition}

\begin{lemma}[The factors of the symbol product]
\label{lem:daub-symbol-product-factors}
\lean{Daubechies.norm_symbol_div_two_pow_sub_one_le, Daubechies.summable_norm_symbol_sub_one,
  Daubechies.multipliable_symbolProduct}
\leanok
\uses{def:daub-symbol-product, lem:daub-symbol-lipschitz}
For every Daubechies mask, every $\xi \in \R$ and every $j \in \N$,
\[
  \Bigl|m_0\Bigl(\frac{\xi}{2^{j+1}}\Bigr) - 1\Bigr| \leq \frac{C_0 |\xi|}{2^{j+1}}
\]
with the constant $C_0$ of \Cref{lem:daub-symbol-lipschitz}. Moreover
$\sum_{j \geq 0} |m_0(\xi/2^{j+1}) - 1| < \infty$, and the product $\Phi(\xi)$ converges.
\end{lemma}

\begin{proof}
\leanok
\uses{lem:daub-symbol-lipschitz}
\Cref{lem:daub-symbol-lipschitz} at the point $\xi/2^{j+1}$. The bounds form a geometric
series, and a product $\prod_j (1 + a_j)$ of complex numbers with $\sum_j |a_j| < \infty$
converges.
\end{proof}

\begin{theorem}[Locally uniform convergence]
\label{thm:daub-symbol-product-continuous}
\lean{Daubechies.hasProdUniformlyOn_symbolProduct,
  Daubechies.hasProdLocallyUniformly_symbolProduct,
  Daubechies.tendstoLocallyUniformly_prod_symbol, Daubechies.tendsto_prod_symbol,
  Daubechies.continuous_symbolProduct}
\leanok
\uses{def:daub-symbol-product}
For every Daubechies mask the product $\Phi$ converges uniformly on every compact set
$K \subseteq \R$. The partial products $\prod_{j<n} m_0(\xi/2^{j+1})$ converge to $\Phi(\xi)$ as
$n \to \infty$, for every $\xi \in \R$ and locally uniformly on $\R$, and $\Phi$ is continuous.
\end{theorem}

\begin{proof}
\leanok
\uses{lem:daub-symbol-product-factors, lem:daub-symbol-props}
On a compact set $K \subseteq [-R, R]$ the bound of \Cref{lem:daub-symbol-product-factors} is at
most $C_0 R/2^{j+1}$, uniformly in $\xi$, and the factors are continuous
(\Cref{lem:daub-symbol-props}); Mathlib's uniform version of the convergence criterion for
products $\prod_j (1 + a_j)$ applies. Locally uniform limits of continuous functions are
continuous.
\end{proof}

\paragraph{Algebraic properties.}
The product inherits the bounds and the symmetry of the symbol, and splitting off its first
factor gives the refinement identity.

\begin{proposition}[Values, bounds and symmetry of $\Phi$]
\label{prop:daub-symbol-product-basic}
\lean{Daubechies.symbolProduct_zero, Daubechies.norm_prod_symbol_le_one,
  Daubechies.norm_symbolProduct_le_one, Daubechies.symbolProduct_neg}
\leanok
\uses{def:daub-symbol-product}
For every Daubechies mask, $\Phi(0) = 1$, and for every $\xi \in \R$ and every $n \in \N$,
\[
  \Bigl|\prod_{j<n} m_0\Bigl(\frac{\xi}{2^{j+1}}\Bigr)\Bigr| \leq 1 , \qquad
  |\Phi(\xi)| \leq 1 , \qquad \Phi(-\xi) = \overline{\Phi(\xi)} .
\]
\end{proposition}

\begin{proof}
\leanok
\uses{lem:daub-symbol-sum, thm:daub-qmf, lem:daub-symbol-props,
  lem:daub-symbol-product-factors}
$m_0(0) = 1$ (\Cref{lem:daub-symbol-sum}) and $|m_0| \leq 1$ (\Cref{thm:daub-qmf}), and the
modulus of a convergent product is the product of the moduli
(\Cref{lem:daub-symbol-product-factors}). Conjugation is continuous and multiplicative, and
$m_0(-\eta) = \overline{m_0(\eta)}$ (\Cref{lem:daub-symbol-props}).
\end{proof}

\begin{lemma}[Splitting off factors]
\label{lem:daub-symbol-product-split}
\lean{Daubechies.prod_range_succ_symbol, Daubechies.symbolProduct_eq_prod_mul}
\leanok
\uses{def:daub-symbol-product}
For every Daubechies mask, every $\xi \in \R$ and every $n \in \N$,
\[
  \prod_{j<n+1} m_0\Bigl(\frac{\xi}{2^{j+1}}\Bigr)
  = m_0\Bigl(\frac{\xi}{2}\Bigr) \prod_{j<n} m_0\Bigl(\frac{\xi/2}{2^{j+1}}\Bigr)
\]
and
\[
  \Phi(\xi) = \prod_{j<n} m_0\Bigl(\frac{\xi}{2^{j+1}}\Bigr) \cdot \Phi\Bigl(\frac{\xi}{2^n}\Bigr) .
\]
\end{lemma}

\begin{proof}
\leanok
\uses{lem:daub-dyadic, lem:daub-symbol-product-factors}
$\xi/2^{j+2} = (\xi/2)/2^{j+1}$ and $\xi/2^{n+j+1} = (\xi/2^n)/2^{j+1}$
(\Cref{lem:daub-dyadic}), and a convergent product splits into its first $n$ factors and its
tail, which converges (\Cref{lem:daub-symbol-product-factors}).
\end{proof}

\begin{theorem}[The refinement identity]
\label{thm:daub-refinement-identity}
\lean{Daubechies.symbolProduct_two_mul}
\leanok
\uses{def:daub-symbol-product}
For every Daubechies mask and every $\xi \in \R$, $\Phi(2\xi) = m_0(\xi) \Phi(\xi)$.
\end{theorem}

\begin{proof}
\leanok
\uses{lem:daub-symbol-product-split}
\Cref{lem:daub-symbol-product-split} with $n = 1$ at the point $2\xi$.
\end{proof}

\begin{proposition}[$\Phi$ near the origin]
\label{prop:daub-symbol-product-nonzero}
\lean{Daubechies.symbolProduct_ne_zero, Daubechies.exists_pos_le_norm_symbolProduct}
\leanok
\uses{def:daub-symbol-product}
For every Daubechies mask there is $\kappa > 0$ such that $\Phi(\xi) \neq 0$ and
$|\Phi(\xi)| \geq \kappa$ for every $\xi \in \R$ with $|\xi| \leq \frac12$.
\end{proposition}

\begin{proof}
\leanok
\uses{lem:daub-symbol-props, lem:daub-symbol-product-factors,
  thm:daub-symbol-product-continuous}
For $|\xi| \leq \frac12$ every argument $\xi/2^{j+1}$ lies in $[-\frac14, \frac14]$, where
$m_0 \neq 0$ (\Cref{lem:daub-symbol-props}), and a convergent product of nonzero factors
$1 + a_j$ with $\sum_j |a_j| < \infty$ is nonzero (\Cref{lem:daub-symbol-product-factors}). The
continuous function $|\Phi|$ (\Cref{thm:daub-symbol-product-continuous}) attains its minimum on
the compact interval $[-\frac12, \frac12]$.
\end{proof}

\paragraph{The modulus.}
The squared modulus of each factor is $\cos(\pi\eta)^{2N} P(\eta)$, and the cosines multiply to
the function $\operatorname{sinc}$.

\begin{theorem}[The modulus of the symbol product]
\label{thm:daub-symbol-product-modulus}
\lean{Daubechies.normSq_symbolProduct}
\leanok
\uses{def:daub-symbol-product, def:daub-fun}
For every Daubechies mask of order $N$ and every $\xi \in \R$,
\[
  |\Phi(\xi)|^2
  = \operatorname{sinc}(\pi\xi)^{2N} \prod_{j \geq 0} P\Bigl(\frac{\xi}{2^{j+1}}\Bigr) .
\]
\end{theorem}

\begin{proof}
\leanok
\uses{prop:daub-symbol-modulus, thm:daub-cos-product, lem:daub-fun-product,
  lem:daub-symbol-product-factors}
The factors satisfy $|m_0(\eta)|^2 = \cos(\pi\eta)^{2N} P(\eta)$
(\Cref{prop:daub-symbol-modulus}), the cosines multiply to
$\prod_{j \geq 0} \cos(\pi\xi/2^{j+1}) = \operatorname{sinc}(\pi\xi)$
(\Cref{thm:daub-cos-product}), and the products of $P$ and of the factors of $\Phi$ converge
(\Cref{lem:daub-fun-product,lem:daub-symbol-product-factors}). The factorwise product of two
convergent products converges to the product of their values.
\end{proof}

\begin{proposition}[A bound by finite products]
\label{prop:daub-symbol-product-bound}
\lean{Daubechies.normSq_symbolProduct_le}
\leanok
\uses{def:daub-symbol-product, def:daub-fun}
For every Daubechies mask of order $N$, every $n \in \N$ and every $\xi \in \R$ with
$|\xi| \leq 2^n/2$,
\[
  |\Phi(\xi)|^2 \leq \exp\Bigl(\binom{2N-1}{N} \frac{\pi^2}{12}\Bigr)
  \operatorname{sinc}(\pi\xi)^{2N} \prod_{j<n} P\Bigl(\frac{\xi}{2^{j+1}}\Bigr) .
\]
\end{proposition}

\begin{proof}
\leanok
\uses{thm:daub-symbol-product-modulus, lem:daub-fun-product, prop:daub-fun-product-bound}
Split the product of $P$ in \Cref{thm:daub-symbol-product-modulus} into its first $n$ factors
and its tail (\Cref{lem:daub-fun-product}), and bound the tail by
\Cref{prop:daub-fun-product-bound}.
\end{proof}

The next section bounds the integrals of $\prod_{j<n} P(\xi/2^{j+1})$ through the transfer
operator.

\section{The transfer operator}

The transfer operator of Eirola~\cite{Eirola} and Villemoes~\cite{Villemoes} turns the decay of
the symbol product into the growth of its iterates: a bound $T^n 1 \leq C \lambda^n$ gives
$\int_\R (1 + |\xi|)^{2s} |\Phi(\xi)|^2 \, d\xi < \infty$ for every $s < N - \log_4 \lambda$. For
$N = 2$ and $N = 3$ explicit eigenfunctions give $\lambda = 4$ and $\lambda = 9$, hence the
exponents $1$ and $3 - \log_2 3$.

\paragraph{The operator.}
The points $\xi/2$ and $(\xi + 1)/2$ are the two preimages of $\xi$ under the doubling map
$\eta \mapsto 2\eta$ modulo $1$.

\begin{definition}[The transfer operator]
\label{def:daub-transfer}
\lean{Daubechies.transfer, Daubechies.transfer_add, Daubechies.transfer_smul}
\leanok
\uses{def:daub-fun}
The \emph{transfer operator} $T$ of order $N \in \N$ maps a function $g : \R \to \R$ to
\[
  (T g)(\xi) = P\Bigl(\frac{\xi}{2}\Bigr) g\Bigl(\frac{\xi}{2}\Bigr)
    + P\Bigl(\frac{\xi + 1}{2}\Bigr) g\Bigl(\frac{\xi + 1}{2}\Bigr) , \qquad \xi \in \R .
\]
It is linear.
\end{definition}

\begin{lemma}[Elementary properties of $T$]
\label{lem:daub-transfer-basic}
\lean{Daubechies.transfer_periodic, Daubechies.continuous_transfer, Daubechies.transfer_nonneg,
  Daubechies.transfer_mono, Daubechies.abs_transfer_le}
\leanok
\uses{def:daub-transfer}
Let $N \in \N$, and let $g$, $g_1$ and $g_2$ be real functions on $\R$. If $g$ is continuous,
so is $T g$; if $g$ has period $1$, so has $T g$; if $g \geq 0$, then $T g \geq 0$; and if
$g_1 \leq g_2$, then $T g_1 \leq T g_2$. If $|g| \leq B$ on $\R$, then
$|T g| \leq 2 \binom{2N-1}{N} B$ on $\R$.
\end{lemma}

\begin{proof}
\leanok
\uses{lem:daub-fun-symmetry, prop:daub-fun-bounds}
$P$ is continuous, has period $1$ (\Cref{lem:daub-fun-symmetry}) and satisfies
$0 \leq P \leq \binom{2N-1}{N}$ (\Cref{prop:daub-fun-bounds}). From $\xi$ to $\xi + 1$ the points
$\xi/2$ and $(\xi + 1)/2$ become $(\xi + 1)/2$ and $\xi/2 + 1$, so for $1$-periodic $g$ the two
summands swap.
\end{proof}

\begin{lemma}[The iterates of $T$]
\label{lem:daub-transfer-iterates}
\lean{Daubechies.continuous_iterate_transfer, Daubechies.periodic_iterate_transfer,
  Daubechies.iterate_transfer_nonneg, Daubechies.abs_iterate_transfer_le,
  Daubechies.continuous_iterate_transfer_one, Daubechies.periodic_iterate_transfer_one,
  Daubechies.iterate_transfer_one_nonneg}
\leanok
\uses{def:daub-transfer}
Let $N, n \in \N$ and $g : \R \to \R$. If $g$ is continuous, has period $1$, or is nonnegative,
then so is $T^n g$. If $|g| \leq B$ on $\R$, then $|T^n g| \leq (2 \binom{2N-1}{N})^n B$ on $\R$.
In particular $T^n 1$ is continuous, nonnegative and has period $1$.
\end{lemma}

\begin{proof}
\leanok
\uses{lem:daub-transfer-basic}
Induction on $n$ with \Cref{lem:daub-transfer-basic}.
\end{proof}

\paragraph{The adjoint identity.}
Integrated against a $1$-periodic function $h$, the operator $T$ becomes the multiplication by
$P$ together with the dilation of $h$ by $2$. Iterated, this identifies the integrals of the
finite products of $P$ with the integrals of $T^n 1$.

\begin{lemma}[The two branches of halving]
\label{lem:daub-halving}
\lean{Daubechies.periodic_comp_two_mul,
  Daubechies.integral_comp_div_two_add_comp_add_one_div_two}
\leanok
If $h : \R \to \R$ has period $1$, then so has $\eta \mapsto h(2\eta)$. For every continuous
$f : \R \to \R$ with period $1$,
\[
  \int_{-1/2}^{1/2} \Bigl(f\Bigl(\frac{\eta}{2}\Bigr) + f\Bigl(\frac{\eta + 1}{2}\Bigr)\Bigr)
  \, d\eta = 2 \int_{-1/2}^{1/2} f(\eta) \, d\eta .
\]
\end{lemma}

\begin{proof}
\leanok
The substitutions $u = \eta/2$ and $u = (\eta + 1)/2$ give $2 \int_{-1/4}^{1/4} f$ and
$2 \int_{1/4}^{3/4} f$, together $2 \int_{-1/4}^{3/4} f$, and the integral of a $1$-periodic
function over an interval of length $1$ does not depend on the interval.
\end{proof}

\begin{theorem}[The adjoint identity]
\label{thm:daub-transfer-adjoint}
\lean{Daubechies.integral_transfer}
\leanok
\uses{def:daub-transfer}
For every $N \in \N$ and all continuous $h, g : \R \to \R$ with period $1$,
\[
  \int_{-1/2}^{1/2} h(\eta) \, (T g)(\eta) \, d\eta
  = 2 \int_{-1/2}^{1/2} h(2\eta) P(\eta) g(\eta) \, d\eta .
\]
\end{theorem}

\begin{proof}
\leanok
\uses{lem:daub-halving, lem:daub-fun-symmetry}
The function $f(u) = h(2u) P(u) g(u)$ is continuous and has period $1$
(\Cref{lem:daub-halving,lem:daub-fun-symmetry}), and the integrand on the left is
$f(\eta/2) + f((\eta + 1)/2)$, since $h(\eta + 1) = h(\eta)$. Apply \Cref{lem:daub-halving}.
\end{proof}

\begin{proposition}[The adjoint identity for the iterates]
\label{prop:daub-transfer-iterates-adjoint}
\lean{Daubechies.integral_iterate_transfer, Daubechies.integral_iterate_transfer_one}
\leanok
\uses{def:daub-transfer}
For all $N, n \in \N$ and all continuous $h, g : \R \to \R$ with period $1$,
\[
  \int_{-1/2}^{1/2} h(\eta) \, (T^n g)(\eta) \, d\eta
  = 2^n \int_{-1/2}^{1/2} h(2^n \eta) \prod_{i<n} P(2^i \eta) \, g(\eta) \, d\eta .
\]
In particular
\[
  \int_{-1/2}^{1/2} (T^n 1)(\eta) \, d\eta
  = 2^n \int_{-1/2}^{1/2} \prod_{i<n} P(2^i \eta) \, d\eta .
\]
\end{proposition}

\begin{proof}
\leanok
\uses{thm:daub-transfer-adjoint, lem:daub-transfer-iterates, lem:daub-halving}
Induction on $n$: $T^{n+1} g = T (T^n g)$, \Cref{thm:daub-transfer-adjoint} for the continuous
$1$-periodic function $T^n g$ (\Cref{lem:daub-transfer-iterates}), and the induction hypothesis
for the continuous $1$-periodic function $\eta \mapsto h(2\eta) P(\eta)$ in place of $h$
(\Cref{lem:daub-halving}).
\end{proof}

\begin{theorem}[The finite products and the transfer operator]
\label{thm:daub-transfer-products}
\lean{Daubechies.prod_daubFun_two_pow_mul_div, Daubechies.integral_prod_daubFun}
\leanok
\uses{def:daub-transfer}
For all $N, n \in \N$,
\[
  \int_{-2^n/2}^{2^n/2} \prod_{j<n} P\Bigl(\frac{\xi}{2^{j+1}}\Bigr) \, d\xi
  = \int_{-1/2}^{1/2} (T^n 1)(\eta) \, d\eta .
\]
\end{theorem}

\begin{proof}
\leanok
\uses{prop:daub-transfer-iterates-adjoint}
The substitution $\xi = 2^n \eta$ turns the left side into
\[
  2^n \int_{-1/2}^{1/2} \prod_{j<n} P\Bigl(\frac{2^n \eta}{2^{j+1}}\Bigr) \, d\eta ,
\]
and $\prod_{j<n} P(2^n \eta/2^{j+1}) = \prod_{i<n} P(2^i \eta)$ for every $\eta \in \R$, by the
reindexing $i = n - 1 - j$. Apply \Cref{prop:daub-transfer-iterates-adjoint}.
\end{proof}

\paragraph{Growth bounds.}
Bounds on $T^n 1$ thereby bound the integrals of the finite products of $P$.

\begin{definition}[Growth bounds]
\label{def:daub-transfer-growth}
\lean{Daubechies.TransferGrowth, Daubechies.TransferGrowth.one_le}
\leanok
\uses{def:daub-transfer}
Let $N \in \N$. A \emph{growth bound} for the transfer operator of order $N$ is a pair of real
numbers $\lambda > 0$ and $C$ with
\[
  (T^n 1)(\xi) \leq C \lambda^n \quad\text{for every } n \in \N \text{ and every } \xi \in \R .
\]
Every growth bound has $C \geq 1$.
\end{definition}

\begin{lemma}[Integrals under a growth bound]
\label{lem:daub-growth-integrals}
\lean{Daubechies.prod_daubFun_nonneg, Daubechies.continuous_prod_daubFun,
  Daubechies.TransferGrowth.integral_iterate_transfer_one_le,
  Daubechies.TransferGrowth.integral_prod_daubFun_le,
  Daubechies.TransferGrowth.setIntegral_prod_daubFun_le}
\leanok
\uses{def:daub-transfer-growth}
Let $N, n \in \N$. The function $\xi \mapsto \prod_{j<n} P(\xi/2^{j+1})$ is continuous and
nonnegative. If $\lambda, C$ is a growth bound for the transfer operator of order $N$, then for
every set $S \subseteq [-2^n/2, 2^n/2]$
\[
  \int_{-1/2}^{1/2} (T^n 1)(\eta) \, d\eta \leq C \lambda^n
  \quad\text{and}\quad
  \int_S \prod_{j<n} P\Bigl(\frac{\xi}{2^{j+1}}\Bigr) \, d\xi \leq C \lambda^n ,
\]
in particular for $S = [-2^n/2, 2^n/2]$.
\end{lemma}

\begin{proof}
\leanok
\uses{thm:daub-transfer-products, lem:daub-fun-symmetry, prop:daub-fun-bounds}
$P$ is continuous and nonnegative (\Cref{lem:daub-fun-symmetry,prop:daub-fun-bounds}). The
first bound integrates $T^n 1 \leq C \lambda^n$ over an interval of length $1$. The integrand of
the second is nonnegative, so the integral over $S$ is at most the integral over
$[-2^n/2, 2^n/2]$, which equals the first integral by \Cref{thm:daub-transfer-products}.
\end{proof}

\begin{proposition}[The crude growth bound]
\label{prop:daub-crude-growth}
\lean{Daubechies.transferGrowth_crude, Daubechies.logb_four_two_mul_choose_lt}
\leanok
\uses{def:daub-transfer-growth}
For every $N \in \N$, $\lambda = 2 \binom{2N-1}{N}$ and $C = 1$ form a growth bound for the
transfer operator of order $N$. For every $N \geq 2$,
$\log_4 \bigl(2 \binom{2N-1}{N}\bigr) < N - \frac12$.
\end{proposition}

\begin{proof}
\leanok
\uses{lem:daub-transfer-iterates, lem:daub-choose-bound}
\Cref{lem:daub-transfer-iterates} with $g = 1$ and $B = 1$. For $N \geq 2$,
$2 \binom{2N-1}{N} < 2 \cdot 4^{N-1} = 4^{N - 1/2}$ by \Cref{lem:daub-choose-bound}.
\end{proof}

\paragraph{Decay of the symbol product.}
Under a growth bound the integrals of $|\Phi|^2$ over the dyadic annuli decay geometrically,
since $\operatorname{sinc}(\pi\xi)^{2N}$ is at most $(\pi |\xi|)^{-2N}$.

\begin{definition}[Dyadic annuli]
\label{def:daub-annulus}
\lean{Daubechies.annulus, Daubechies.measurableSet_annulus, Daubechies.annulus_subset_Icc,
  Daubechies.exists_mem_annulus, Daubechies.univ_subset_Icc_union_annulus}
\leanok
For $n \in \N$ the \emph{dyadic annulus} is
\[
  A_n = \{\xi \in \R : 2^n/4 < |\xi| \leq 2^n/2\} .
\]
It is measurable and contained in $[-2^n/2, 2^n/2]$. Every $\xi \in \R$ with $|\xi| > \frac14$
lies in some $A_n$, so $\R = [-\frac14, \frac14] \cup \bigcup_{n \in \N} A_n$.
\end{definition}

\begin{lemma}[Integrability from the annuli]
\label{lem:daub-annulus-integrable}
\lean{Daubechies.integrable_of_integral_annulus_le}
\leanok
\uses{def:daub-annulus}
Let $f$ be a continuous function from $\R$ to a normed abelian group, and let $A \in \R$ and
$0 \leq r < 1$ with $\int_{A_n} |f| \leq A r^n$ for every $n \in \N$. Then $f$ is integrable on
$\R$.
\end{lemma}

\begin{proof}
\leanok
On $[-\frac14, \frac14]$ the function $f$ is integrable as a continuous function on a compact
interval, and on $\bigcup_n A_n$ since the integrals over the annuli are summable; these two
sets cover $\R$ (\Cref{def:daub-annulus}).
\end{proof}

\begin{lemma}[$\Phi$ on the annuli]
\label{lem:daub-sinc-annulus}
\lean{Daubechies.abs_sinc_le_inv_abs, Daubechies.sinc_pow_le_of_mem_annulus,
  Daubechies.normSq_symbolProduct_le_of_mem_annulus}
\leanok
\uses{def:daub-annulus, def:daub-symbol-product}
For every $x \in \R$ with $x \neq 0$, $|\operatorname{sinc} x| \leq 1/|x|$. For all
$N, n \in \N$ and every $\xi \in A_n$,
$\operatorname{sinc}(\pi\xi)^{2N} \leq (4/(\pi 2^n))^{2N}$, and for every Daubechies mask of
order $N$
\[
  |\Phi(\xi)|^2 \leq \exp\Bigl(\binom{2N-1}{N} \frac{\pi^2}{12}\Bigr)
  \Bigl(\frac{4}{\pi 2^n}\Bigr)^{2N} \prod_{j<n} P\Bigl(\frac{\xi}{2^{j+1}}\Bigr) .
\]
\end{lemma}

\begin{proof}
\leanok
\uses{prop:daub-symbol-product-bound, lem:daub-growth-integrals}
$|\sin x| \leq 1$, and $\pi |\xi| > \pi 2^n/4$ on $A_n$. Since $|\xi| \leq 2^n/2$ on $A_n$,
\Cref{prop:daub-symbol-product-bound} applies, and the finite product is nonnegative
(\Cref{lem:daub-growth-integrals}).
\end{proof}

\begin{proposition}[$|\Phi|^2$ on the annuli]
\label{prop:daub-annulus-integral}
\lean{Daubechies.integral_annulus_normSq_le}
\leanok
\uses{def:daub-annulus, def:daub-symbol-product, def:daub-transfer-growth}
Let $\lambda, C$ be a growth bound for the transfer operator of order $N$. Then for every
Daubechies mask of order $N$ and every $n \in \N$,
\[
  \int_{A_n} |\Phi(\xi)|^2 \, d\xi
  \leq \exp\Bigl(\binom{2N-1}{N} \frac{\pi^2}{12}\Bigr)
  \Bigl(\frac{4}{\pi 2^n}\Bigr)^{2N} C \lambda^n .
\]
\end{proposition}

\begin{proof}
\leanok
\uses{lem:daub-sinc-annulus, lem:daub-growth-integrals, thm:daub-symbol-product-continuous}
Both sides of the bound of \Cref{lem:daub-sinc-annulus} are continuous
(\Cref{thm:daub-symbol-product-continuous,lem:daub-growth-integrals}); integrate it over
$A_n \subseteq [-2^n/2, 2^n/2]$ and apply \Cref{lem:daub-growth-integrals}.
\end{proof}

Since $(4/(\pi 2^n))^{2N} = (4/\pi)^{2N} 4^{-Nn}$, these bounds decay geometrically in $n$, with
the ratio $\lambda/4^N$.

\begin{lemma}[Weights on the annuli]
\label{lem:daub-weights}
\lean{Daubechies.rpow_le_four_rpow_abs_mul, Daubechies.one_add_abs_rpow_le_of_mem_annulus,
  Daubechies.four_div_pi_mul_two_pow_pow, Daubechies.two_rpow_two_mul_mul_lt}
\leanok
\uses{def:daub-annulus}
The following hold.
\begin{enumerate}
  \item If $x, y \in \R$ satisfy $y > 0$ and $y/4 \leq x \leq 4y$, then $x^t \leq 4^{|t|} y^t$ for
    every $t \in \R$.
  \item For every $t \in \R$, every $n \in \N$ and every $\xi \in A_n$,
    $(1 + |\xi|)^t \leq 4^{|t|} (2^t)^n$.
  \item For all $N, n \in \N$, $(4/(\pi 2^n))^{2N} = (4/\pi)^{2N}/(4^N)^n$.
  \item If $N \in \N$ and $\lambda, s \in \R$ satisfy $\lambda > 0$ and
    $s < N - \log_4 \lambda$, then $2^{2s} \lambda < 4^N$.
\end{enumerate}
\end{lemma}

\begin{proof}
\leanok
Item~1 follows from $x \leq 4y$ for $t \geq 0$ and from $y/4 \leq x$ for $t \leq 0$. Item~2 is
item~1 with $y = 2^n$, since $2^n/4 \leq 1 + |\xi| \leq 4 \cdot 2^n$ on $A_n$. Item~3 is
$(2^n)^{2N} = (4^N)^n$, and item~4 says $\lambda < 4^{N-s} = 4^N/2^{2s}$.
\end{proof}

\begin{lemma}[Splitting a product of powers]
\label{lem:daub-weight-split}
\lean{Daubechies.rpow_mul_le_rpow_add_rpow_mul_sq}
\leanok
For every $x > 0$, every $c \geq 0$ and all $a, b \in \R$,
$x^a c \leq x^{2(a-b)} + x^{2b} c^2$.
\end{lemma}

\begin{proof}
\leanok
With $u = x^{a-b}$ and $v = x^b c$, $x^a c = uv \leq 2uv \leq u^2 + v^2$.
\end{proof}

\begin{theorem}[Weighted $L_2$ decay of the symbol product]
\label{thm:daub-decay-l2}
\lean{Daubechies.integrable_rpow_mul_normSq}
\leanok
\uses{def:daub-symbol-product, def:daub-transfer-growth}
Let $\lambda, C$ be a growth bound for the transfer operator of order $N$. Then for every
Daubechies mask of order $N$ and every $s \in \R$ with $s < N - \log_4 \lambda$, the function
$\xi \mapsto (1 + |\xi|)^{2s} |\Phi(\xi)|^2$ is integrable on $\R$.
\end{theorem}

\begin{proof}
\leanok
\uses{lem:daub-weights, prop:daub-annulus-integral, lem:daub-annulus-integrable,
  thm:daub-symbol-product-continuous}
On $A_n$ the weight is at most $4^{|2s|} (2^{2s})^n$ (\Cref{lem:daub-weights}), so by
\Cref{prop:daub-annulus-integral} the integral over $A_n$ is at most a constant times
$(2^{2s} \lambda/4^N)^n$, and $2^{2s} \lambda/4^N < 1$ (\Cref{lem:daub-weights}). The function is
continuous (\Cref{thm:daub-symbol-product-continuous}); apply \Cref{lem:daub-annulus-integrable}.
\end{proof}

\begin{theorem}[Weighted $L_1$ decay of the symbol product]
\label{thm:daub-decay-l1}
\lean{Daubechies.integrable_rpow_mul_norm}
\leanok
\uses{def:daub-symbol-product, def:daub-transfer-growth}
Let $\lambda, C$ be a growth bound for the transfer operator of order $N$. Then for every
Daubechies mask of order $N$ and every $a \in \R$ with $a < N - \log_4 \lambda - \frac12$, the
function $\xi \mapsto (1 + |\xi|)^a |\Phi(\xi)|$ is integrable on $\R$.
\end{theorem}

\begin{proof}
\leanok
\uses{lem:daub-weight-split, thm:daub-decay-l2, thm:daub-symbol-product-continuous}
Choose $b$ with $a + \frac12 < b < N - \log_4 \lambda$. By \Cref{lem:daub-weight-split},
\[
  (1 + |\xi|)^a |\Phi(\xi)| \leq (1 + |\xi|)^{-2(b - a)} + (1 + |\xi|)^{2b} |\Phi(\xi)|^2 ,
\]
where the first summand is integrable since $2(b - a) > 1$, and the second by
\Cref{thm:daub-decay-l2}. The function on the left is continuous
(\Cref{thm:daub-symbol-product-continuous}).
\end{proof}

\begin{corollary}[The symbol product is integrable and square integrable]
\label{cor:daub-symbol-product-integrable}
\lean{Daubechies.integrable_symbolProduct, Daubechies.memLp_symbolProduct_two}
\leanok
\uses{def:daub-symbol-product}
For every Daubechies mask of order $N \geq 2$ the symbol product $\Phi$ is integrable and square
integrable on $\R$.
\end{corollary}

\begin{proof}
\leanok
\uses{prop:daub-crude-growth, thm:daub-decay-l1, thm:daub-decay-l2}
The crude growth bound has $\log_4 \lambda < N - \frac12$ (\Cref{prop:daub-crude-growth}); take
$a = 0$ in \Cref{thm:daub-decay-l1} and $s = 0$ in \Cref{thm:daub-decay-l2}.
\end{proof}

\paragraph{Orders two and three.}
For $N = 2$ and $N = 3$ the transfer operator has eigenfunctions among the trigonometric
polynomials, and $T^n 1$ is known exactly. In the variable $y = \sin^2(\pi\xi/2)$ every
eigenvalue equation becomes a polynomial identity in $y$.

\begin{lemma}[Linearity of the iterates and eigenfunctions]
\label{lem:daub-transfer-eigen}
\lean{Daubechies.iterate_transfer_add, Daubechies.iterate_transfer_smul,
  Daubechies.iterate_transfer_eq_pow_smul}
\leanok
\uses{def:daub-transfer}
For all $N, n \in \N$, all $g, g_1, g_2 : \R \to \R$ and every $c \in \R$,
$T^n (g_1 + g_2) = T^n g_1 + T^n g_2$ and $T^n (c g) = c \, T^n g$. If $e : \R \to \R$ and
$\mu \in \R$ satisfy $T e = \mu e$, then $T^n e = \mu^n e$ for every $n \in \N$.
\end{lemma}

\begin{proof}
\leanok
Induction on $n$, since $T$ is linear (\Cref{def:daub-transfer}).
\end{proof}

\begin{lemma}[The variable $y$]
\label{lem:daub-transfer-sin-sq}
\lean{Daubechies.transfer_eq_eval_sin_sq, Daubechies.cos_two_pi_mul_div_two,
  Daubechies.cos_two_pi_mul_add_one_div_two, Daubechies.cos_two_pi_mul_eq_sin_sq}
\leanok
\uses{def:daub-transfer, def:daub-poly}
Let $\xi \in \R$ and $y = \sin^2(\pi\xi/2)$. For every $N \in \N$ and every $g : \R \to \R$,
\[
  (T g)(\xi) = P_N(y) \, g\Bigl(\frac{\xi}{2}\Bigr)
    + P_N(1 - y) \, g\Bigl(\frac{\xi + 1}{2}\Bigr) ,
\]
and $\cos \pi\xi = 1 - 2y$, $\cos(\pi(\xi + 1)) = 2y - 1$ and $\cos 2\pi\xi = 1 - 8y(1 - y)$.
\end{lemma}

\begin{proof}
\leanok
\uses{lem:daub-fun-symmetry}
$P(\xi/2) = P_N(y)$ and
$P((\xi + 1)/2) = P(\xi/2 + \frac12) = P_N(\cos^2(\pi\xi/2)) = P_N(1 - y)$
(\Cref{lem:daub-fun-symmetry}). The double angle formulas give $\cos \pi\xi = 1 - 2y$ and
$\cos 2\pi\xi = 1 - 2 \sin^2 \pi\xi = 1 - 8y(1 - y)$, and $\cos(\pi\xi + \pi) = -\cos \pi\xi$.
\end{proof}

\begin{proposition}[Order two]
\label{prop:daub-transfer-two}
\lean{Daubechies.transfer_two_one, Daubechies.iterate_transfer_two_one,
  Daubechies.transferGrowth_two, Daubechies.two_sub_logb_four_four}
\leanok
\uses{def:daub-transfer-growth}
For $N = 2$, $T 1 = 4$ and $T^n 1 = 4^n$ for every $n \in \N$. The numbers $\lambda = 4$ and
$C = 1$ form a growth bound for the transfer operator of order $2$, and $2 - \log_4 4 = 1$.
\end{proposition}

\begin{proof}
\leanok
\uses{lem:daub-transfer-sin-sq, lem:daub-transfer-eigen, def:daub-poly}
$P_2(y) = 1 + 2y$ (\Cref{def:daub-poly}), so $(T 1)(\xi) = (1 + 2y) + (1 + 2(1 - y)) = 4$ by
\Cref{lem:daub-transfer-sin-sq}. Then \Cref{lem:daub-transfer-eigen} with $e = 1$ and $\mu = 4$.
\end{proof}

\begin{proposition}[Eigenfunctions of order three]
\label{prop:daub-transfer-three-eigen}
\lean{Daubechies.transfer_three_eigen_one, Daubechies.transfer_three_eigen_two}
\leanok
\uses{def:daub-transfer}
Let $N = 3$ and $c(\xi) = \cos 2\pi\xi$. Then
\[
  T (9 + c) = 9 \, (9 + c) \quad\text{and}\quad T (1 + 3c) = -4 \, (1 + 3c) .
\]
\end{proposition}

\begin{proof}
\leanok
\uses{lem:daub-transfer-sin-sq, def:daub-poly}
$P_3(y) = 1 + 3y + 6y^2$ (\Cref{def:daub-poly}). By \Cref{lem:daub-transfer-sin-sq}, at $\xi$
both sides of the first identity equal $90 - 72 y(1 - y)$ and both sides of the second equal
$-16 + 96 y(1 - y)$, where $y = \sin^2(\pi\xi/2)$.
\end{proof}

\begin{proposition}[Order three]
\label{prop:daub-transfer-three}
\lean{Daubechies.iterate_transfer_three_one, Daubechies.iterate_transfer_three_one_le,
  Daubechies.transferGrowth_three, Daubechies.three_sub_logb_four_nine}
\leanok
\uses{def:daub-transfer-growth}
For $N = 3$, every $n \in \N$ and every $\xi \in \R$,
\[
  (T^n 1)(\xi) = \frac{3}{26} \, 9^n (9 + \cos 2\pi\xi)
    - \frac{1}{26} \, (-4)^n (1 + 3 \cos 2\pi\xi) ,
\]
and $(T^n 1)(\xi) \leq \frac{17}{13} \, 9^n$. The numbers $\lambda = 9$ and $C = 2$ form a
growth bound for the transfer operator of order $3$, and $3 - \log_4 9 = 3 - \log_2 3$.
\end{proposition}

\begin{proof}
\leanok
\uses{prop:daub-transfer-three-eigen, lem:daub-transfer-eigen}
$1 = \frac{3}{26} (9 + c) - \frac{1}{26} (1 + 3c)$; apply \Cref{lem:daub-transfer-eigen} with the
eigenfunctions of \Cref{prop:daub-transfer-three-eigen}. Then $9 + \cos 2\pi\xi \leq 10$ and
$|(-4)^n (1 + 3 \cos 2\pi\xi)| \leq 4 \cdot 9^n$ give
$T^n 1 \leq (\frac{30}{26} + \frac{4}{26}) \, 9^n$. Finally
$\log_4 9 = \log 3^2 / \log 2^2 = \log_2 3$.
\end{proof}

By \Cref{thm:daub-decay-l2}, $\int_\R (1 + |\xi|)^{2s} |\Phi(\xi)|^2 \, d\xi < \infty$ for every
Daubechies mask of order $2$ and every $s < 1$, and for every Daubechies mask of order $3$ and
every $s < 3 - \log_2 3 \approx 1.415$.

\section{Fourier transforms of integrable functions}

This section is general analysis on $\R$, used for the scaling function
$\varphi = \mathcal F^{-1} \Phi$ in the following sections. The Fourier transform of
$f : \R \to \C$ is
\[
  \widehat f(\xi) = \mathcal F f(\xi) = \int_\R f(x) e^{-2\pi i x \xi} \, dx ,
\]
and its inverse transform is $\mathcal F^{-1} f(x) = \int_\R f(\xi) e^{2\pi i x \xi} \, d\xi$.
Lean reads the integral of a function that is not integrable as $0$; with this convention the
statements below without integrability hypotheses hold for every $f$.

\begin{lemma}[Continuity and the uniform bound]
\label{lem:daub-fourier-continuous}
\lean{Real.continuous_fourier_of_integrable, Real.continuous_fourierInv_of_integrable,
  Real.norm_fourier_le_integral_norm, Real.norm_fourierInv_le_integral_norm}
\leanok
For every integrable $f : \R \to \C$ the functions $\widehat f$ and $\mathcal F^{-1} f$ are
continuous. For every $f : \R \to \C$ and all $\xi, x \in \R$,
\[
  |\widehat f(\xi)| \leq \int_\R |f| \quad\text{and}\quad |\mathcal F^{-1} f(x)| \leq \int_\R |f| .
\]
\end{lemma}

\begin{proof}
\leanok
$|e^{-2\pi i x \xi}| = 1$, and dominated convergence with the dominating function $|f|$; the
inverse transform is $\mathcal F^{-1} f(x) = \widehat f(-x)$.
\end{proof}

\begin{lemma}[Linearity]
\label{lem:daub-fourier-linear}
\lean{Real.fourier_const_mul_apply, Real.fourierInv_const_mul_apply, Real.fourier_sum_apply,
  Real.fourierInv_sum_apply}
\leanok
For every $f : \R \to \C$ and every $c \in \C$, $\mathcal F(c f) = c \widehat f$ and
$\mathcal F^{-1}(c f) = c \mathcal F^{-1} f$. For every finite family $(F_i)_{i \in S}$ of
integrable functions $F_i : \R \to \C$,
\[
  \mathcal F\Bigl(\sum_{i \in S} F_i\Bigr) = \sum_{i \in S} \widehat{F_i}
  \quad\text{and}\quad
  \mathcal F^{-1}\Bigl(\sum_{i \in S} F_i\Bigr) = \sum_{i \in S} \mathcal F^{-1} F_i .
\]
\end{lemma}

\begin{proof}
\leanok
Linearity of the integral.
\end{proof}

\begin{lemma}[Translation and modulation]
\label{lem:daub-fourier-translate}
\lean{Real.fourier_comp_add_right_apply, Real.fourier_comp_sub_right_apply,
  Real.fourierInv_comp_add_right_apply, Real.fourier_exp_mul_apply,
  Real.fourierInv_exp_mul_apply}
\leanok
For $b \in \R$ let $e_b(x) = e^{2\pi i b x}$. For every $f : \R \to \C$ and all
$a, b, \xi, x \in \R$:
\begin{enumerate}
  \item $\mathcal F[f(\cdot + a)](\xi) = e^{2\pi i a \xi} \widehat f(\xi)$ and
    $\mathcal F[f(\cdot - a)](\xi) = e^{-2\pi i a \xi} \widehat f(\xi)$;
  \item $\mathcal F^{-1}[f(\cdot + a)](x) = e^{-2\pi i a x} \mathcal F^{-1} f(x)$;
  \item $\mathcal F(e_b f)(\xi) = \widehat f(\xi - b)$ and
    $\mathcal F^{-1}(e_b f)(x) = \mathcal F^{-1} f(x + b)$.
\end{enumerate}
\end{lemma}

\begin{proof}
\leanok
The substitution $y = x + a$, and $e^{-2\pi i x \xi} e^{2\pi i b x} = e^{-2\pi i x (\xi - b)}$.
\end{proof}

\begin{lemma}[Dilation]
\label{lem:daub-fourier-dilate}
\lean{Real.fourier_comp_mul_left_apply_of_ne_zero, Real.fourier_comp_mul_left_apply,
  Real.fourierInv_comp_mul_left_apply, Real.fourier_comp_mul_sub_apply}
\leanok
For every $f : \R \to \C$ and all $b, \xi, x \in \R$:
\begin{enumerate}
  \item $\mathcal F[f(a\,\cdot)](\xi) = |a|^{-1} \widehat f(\xi / a)$ for every $a \neq 0$;
  \item $\mathcal F^{-1}[f(a\,\cdot)](x) = a^{-1} \mathcal F^{-1} f(x / a)$ for every $a > 0$;
  \item $\widehat g(\xi) = a^{-1} e^{-2\pi i b \xi / a} \widehat f(\xi / a)$ for every $a > 0$,
    where $g(x) = f(ax - b)$.
\end{enumerate}
\end{lemma}

\begin{proof}
\leanok
\uses{lem:daub-fourier-translate}
The substitution $y = ax$, with $dx = |a|^{-1} \, dy$. The third identity is the first one for
$f(\cdot - b)$, together with \Cref{lem:daub-fourier-translate}.
\end{proof}

The third identity is the form met by the cascade algorithm
$\eta_{n+1}(x) = \sum_k c_k \eta_n(2x - k)$.

\begin{lemma}[Conjugation]
\label{lem:daub-fourier-conj}
\lean{Real.conj_fourier, Real.conj_fourierInv}
\leanok
For every $f : \R \to \C$ and all $\xi, x \in \R$, with $\bar f(x) = \overline{f(x)}$,
\[
  \overline{\widehat f(\xi)} = \mathcal F^{-1} \bar f(\xi) \quad\text{and}\quad
  \overline{\mathcal F^{-1} f(x)} = \mathcal F \bar f(x) .
\]
\end{lemma}

\begin{proof}
\leanok
$\overline{e^{-2\pi i x \xi}} = e^{2\pi i x \xi}$, and conjugation commutes with the integral.
\end{proof}

\paragraph{The multiplication formula and Parseval's identity.}
The transform moves from one factor of a product to the other. Together with Fourier inversion
this gives Parseval's identity for integrable functions with integrable transforms.

\begin{theorem}[The multiplication formula]
\label{thm:daub-multiplication-formula}
\lean{Real.integral_fourier_mul, Real.integral_fourierInv_mul}
\leanok
For all integrable $f, g : \R \to \C$,
\[
  \int_\R \widehat f(\xi) \, g(\xi) \, d\xi = \int_\R f(x) \, \widehat g(x) \, dx ,
\]
and the same holds with $\mathcal F^{-1}$ in place of $\mathcal F$:
\[
  \int_\R \mathcal F^{-1} f(\xi) \, g(\xi) \, d\xi = \int_\R f(x) \, \mathcal F^{-1} g(x) \, dx .
\]
\end{theorem}

\begin{proof}
\leanok
By Fubini's theorem both sides of the first identity equal the absolutely convergent double
integral $\int_\R \int_\R e^{-2\pi i x \xi} f(x) g(\xi) \, dx \, d\xi$. The second identity is
the same with $e^{2\pi i x \xi}$.
\end{proof}

\begin{theorem}[Parseval's identity]
\label{thm:daub-parseval}
\lean{Real.integral_mul_conj_eq_integral_fourier, Real.integral_norm_sq_fourier}
\leanok
Let $f : \R \to \C$ be integrable, and let $g : \R \to \C$ be continuous and integrable with
integrable $\widehat g$. Then
\[
  \int_\R f(x) \, \overline{g(x)} \, dx = \int_\R \widehat f(\xi) \, \overline{\widehat g(\xi)}
  \, d\xi .
\]
In particular $\int_\R |\widehat f|^2 = \int_\R |f|^2$ for every continuous integrable
$f : \R \to \C$ with integrable $\widehat f$.
\end{theorem}

\begin{proof}
\leanok
\uses{thm:daub-multiplication-formula, lem:daub-fourier-conj}
By \Cref{lem:daub-fourier-conj} and \Cref{thm:daub-multiplication-formula} for $\mathcal F^{-1}$,
\[
  \int_\R \widehat g \, \overline{\widehat f} = \int_\R \widehat g \, \mathcal F^{-1} \bar f
  = \int_\R \mathcal F^{-1} \widehat g \, \bar f ,
\]
and $\mathcal F^{-1} \widehat g = g$ by the Fourier inversion theorem, which holds for
continuous integrable $g$ with integrable $\widehat g$. Conjugation gives the identity. For
$g = f$ it reads $\int |f|^2 = \int |\widehat f|^2$, since $z \bar z = |z|^2$.
\end{proof}

\begin{lemma}[Testing against Schwartz functions]
\label{lem:daub-fourier-schwartz}
\lean{Real.integral_mul_schwartz_eq}
\leanok
For every integrable $f : \R \to \C$ and every Schwartz function $\psi : \R \to \C$,
\[
  \int_\R f(x) \, \psi(x) \, dx = \int_\R \widehat f(\xi) \, \mathcal F^{-1} \psi(\xi) \, d\xi .
\]
\end{lemma}

\begin{proof}
\leanok
\uses{thm:daub-multiplication-formula}
$\mathcal F^{-1} \psi$ is a Schwartz function, hence integrable, and
$\mathcal F \mathcal F^{-1} \psi = \psi$ by Fourier inversion on the Schwartz space. Apply
\Cref{thm:daub-multiplication-formula} to $f$ and $\mathcal F^{-1} \psi$.
\end{proof}

\paragraph{The fundamental lemma.}
A function whose integrals against all test functions supported in an open set vanish is zero
almost everywhere on that set, and everywhere on it if it is continuous.

\begin{proposition}[The fundamental lemma of the calculus of variations]
\label{prop:daub-fundamental-lemma}
\lean{Real.ae_eq_zero_of_integral_mul_schwartz_eq_zero,
  Real.eqOn_zero_of_integral_mul_schwartz_eq_zero}
\leanok
Let $U \subseteq \R$ be open, and let $f : \R \to \C$ satisfy $\int_\R \psi f = 0$ for every
Schwartz function $\psi : \R \to \C$ whose support is compact and contained in $U$.
\begin{enumerate}
  \item If $f$ is locally integrable on $U$, then $f(x) = 0$ for almost every $x \in U$.
  \item If $f$ is continuous, then $f(x) = 0$ for every $x \in U$.
\end{enumerate}
\end{proposition}

\begin{proof}
\leanok
Mathlib's fundamental lemma tests with real smooth functions $g$ with compact support in $U$,
and $x \mapsto g(x)$, read as a complex function, is a Schwartz function. Two continuous
functions that agree almost everywhere on an open set agree on it, since nonempty open sets
have positive Lebesgue measure.
\end{proof}

\section{Masks of every order}

The masks of every order come from the Fej\'er--Riesz theorem, as in~\cite{Daubechies}: a real
polynomial $Q$ with $Q > 0$ on $[-1, 1]$ satisfies $Q(\cos\theta) = |L(e^{i\theta})|^2$ with a
real polynomial $L$. For $Q(t) = P_N((1 - t)/2)$, where $Q(\cos 2\pi\xi) = P_N(\sin^2 \pi\xi)$,
this gives the polynomial $L$ of a mask of order $N$. The proof factors $Q$ over $\C$ and
splits each factor $\cos\theta - t$ with the inverse of the Joukowsky map.

\paragraph{The inverse of the Joukowsky map.}
The Joukowsky map $a \mapsto (a + a^{-1})/2$ takes $e^{i\varphi}$ to $\cos\varphi$, and for
$a \neq 0$ the equation $(a + a^{-1})/2 = t$ reads $a^2 - 2ta + 1 = 0$.

\begin{lemma}[Roots of $a^2 - 2ta + 1$]
\label{lem:daub-joukowsky-equation}
\lean{Complex.joukowsky_eq_re_of_norm_eq_one, Complex.joukowsky_root_ne_zero,
  Complex.joukowsky_root_conj}
\leanok
Let $a, t \in \C$ with $a^2 - 2ta + 1 = 0$. Then $a \neq 0$ and
$\bar a^2 - 2 \bar t \bar a + 1 = 0$. If $|a| = 1$, then $t = \operatorname{Re} a$; in
particular $t$ is a real number in $[-1, 1]$.
\end{lemma}

\begin{proof}
\leanok
The constant term of the equation is $1$, and conjugation is a ring homomorphism of $\C$. If
$|a| = 1$, then $a^{-1} = \bar a$, and the equation divided by $a$ says
$2t = a + a^{-1} = a + \bar a = 2 \operatorname{Re} a$.
\end{proof}

\begin{lemma}[The root in the disk]
\label{lem:daub-joukowsky-root}
\lean{Complex.exists_joukowsky_root, Complex.joukowsky_root_unique}
\leanok
For every $t \in \C \setminus [-1, 1]$ the equation $a^2 - 2ta + 1 = 0$ has a root $a$ with
$|a| < 1$. For every $t \in \C$ it has at most one root $a$ with $|a| < 1$.
\end{lemma}

\begin{proof}
\leanok
\uses{lem:daub-joukowsky-equation}
The two roots $t \pm s$, $s^2 = t^2 - 1$, have product $1$, and for $t \notin [-1, 1]$ neither
lies on the unit circle (\Cref{lem:daub-joukowsky-equation}), so one of them lies in the open
unit disk. Two different roots are the two roots of the equation and have product $1$, so they
cannot both lie in the open unit disk.
\end{proof}

\begin{definition}[The inverse of the Joukowsky map]
\label{def:daub-joukowsky-inv}
\lean{Complex.joukowskyInv, Complex.joukowskyInv_spec, Complex.joukowskyInv_eq}
\leanok
\uses{lem:daub-joukowsky-root}
For $t \in \C$ let $a(t)$ be the root $a$ of $a^2 - 2ta + 1 = 0$ with $|a| < 1$ if there is
one, and $a(t) = 0$ otherwise. If the equation has a root $a$ with $|a| < 1$, then $a = a(t)$;
in particular $|a(t)| < 1$ and $a(t)^2 - 2t \, a(t) + 1 = 0$. This is the case for every
$t \in \C \setminus [-1, 1]$; on this set $t \mapsto a(t)$ is the inverse of the Joukowsky map
on the punctured open unit disk.
\end{definition}

For $t \in [-1, 1]$ the two roots $t \pm i \sqrt{1 - t^2}$ lie on the unit circle, and
$a(t) = 0$.

\begin{lemma}[The inverse of the Joukowsky map and conjugation]
\label{lem:daub-joukowsky-conj}
\lean{Complex.joukowskyInv_conj}
\leanok
\uses{def:daub-joukowsky-inv}
For every $t \in \C$, $a(\bar t) = \overline{a(t)}$.
\end{lemma}

\begin{proof}
\leanok
\uses{lem:daub-joukowsky-equation, lem:daub-joukowsky-root}
Conjugation maps the roots of $a^2 - 2ta + 1$ in the open unit disk to the roots of
$a^2 - 2\bar t a + 1$ in the open unit disk (\Cref{lem:daub-joukowsky-equation}), and such a root
is unique (\Cref{lem:daub-joukowsky-root}). If there is no such root for $t$, there is none
for $\bar t$, and both sides are $0$.
\end{proof}

\begin{lemma}[The factorisation on the unit circle]
\label{lem:daub-joukowsky-factor}
\lean{Complex.joukowsky_two_mul_mul_re_sub, Complex.joukowsky_two_mul_norm_mul_norm_re_sub}
\leanok
Let $a, t, z \in \C$ with $a^2 - 2ta + 1 = 0$. If $z \bar z = 1$, then
\[
  2a (\operatorname{Re} z - t) = -(1 - az)(1 - a \bar z) ,
\]
and hence $2 |a| \, |\operatorname{Re} z - t| = |1 - az| \, |1 - \bar a z|$ if $|z| = 1$.
\end{lemma}

\begin{proof}
\leanok
Expand the right side with $z \bar z = 1$, $z + \bar z = 2 \operatorname{Re} z$ and
$a^2 + 1 = 2ta$. For the moduli, $|1 - a \bar z| = |\overline{1 - a \bar z}| = |1 - \bar a z|$.
\end{proof}

For $z = e^{i\theta}$ the left side is $2 |a| \, |\cos\theta - t|$: on the unit circle the factor
$\cos\theta - t$ splits into two linear factors in $e^{i\theta}$.

\paragraph{The complex roots of a real polynomial.}
For a real polynomial $Q$, products $\prod_t$ run over the complex roots $t$ of $Q$, counted
with multiplicity.

\begin{lemma}[Real polynomials over $\C$]
\label{lem:daub-real-poly-roots}
\lean{Multiset.norm_prod, Complex.aeval_ofReal, Polynomial.map_conj_aroots,
  Polynomial.map_comp_conj_aroots, Polynomial.ofReal_ne_of_mem_aroots,
  Polynomial.aeval_eq_mul_prod_aroots, Polynomial.eval_eq_abs_leadingCoeff_mul_prod_norm}
\leanok
Let $Q$ be a real polynomial with leading coefficient $c$.
\begin{enumerate}
  \item The conjugates $\bar t$ of the complex roots $t$ of $Q$ are again the complex roots of
    $Q$, with the same multiplicities; so for every function $g$ on $\C$ the values $g(\bar t)$
    and $g(t)$ over the roots agree as multisets.
  \item $Q(z) = c \prod_t (z - t)$ for every $z \in \C$.
  \item $Q(x) = |c| \prod_t |x - t|$ for every $x \in \R$ with $Q(x) \geq 0$.
  \item No $r \in \R$ with $Q(r) \neq 0$ is a complex root of $Q$.
\end{enumerate}
\end{lemma}

\begin{proof}
\leanok
$\C$ is algebraically closed, so $Q$ has $\deg Q$ complex roots, counted with multiplicity, and
item~2 holds; the coefficients of $Q$ are real, so $\bar t$ is a root of the same multiplicity
as $t$. Item~3 is the modulus of item~2, since $Q(x) = |Q(x)|$: the modulus of a product is the
product of the moduli, and a real polynomial takes real values at real points. Item~4 holds
since $Q(t) = 0$ at every root $t$.
\end{proof}

\paragraph{The Fej\'er--Riesz factor.}
For $Q > 0$ on $[-1, 1]$ each root $t$ of $Q$ contributes the factor $1 - a(t) z$, whose zero
$1/a(t)$ lies outside the closed unit disk.

\begin{definition}[The Fej\'er--Riesz factor]
\label{def:daub-fejer-riesz-factor}
\lean{Polynomial.fejerRieszFactor, Polynomial.eval_fejerRieszFactor,
  Polynomial.natDegree_fejerRieszFactor_le}
\leanok
\uses{def:daub-joukowsky-inv}
The \emph{Fej\'er--Riesz factor} of a real polynomial $Q$ is the complex polynomial
\[
  L_0(z) = \prod_t \bigl(1 - a(t) z\bigr) ,
\]
the product over the complex roots $t$ of $Q$, counted with multiplicity. Its degree is at most
$\deg Q$.
\end{definition}

\begin{lemma}[The Fej\'er--Riesz factor is real]
\label{lem:daub-fejer-riesz-factor-real}
\lean{Polynomial.map_conj_fejerRieszFactor, Polynomial.exists_map_eq_of_map_conj_eq}
\leanok
\uses{def:daub-fejer-riesz-factor}
A complex polynomial whose coefficients are fixed by conjugation has real coefficients. For
every real polynomial $Q$ the coefficients of $L_0$ are fixed by conjugation, so $L_0$ is a real
polynomial.
\end{lemma}

\begin{proof}
\leanok
\uses{lem:daub-joukowsky-conj, lem:daub-real-poly-roots}
A complex number fixed by conjugation equals its real part. Conjugating the coefficients of
$L_0$ replaces $a(t)$ by $\overline{a(t)} = a(\bar t)$ (\Cref{lem:daub-joukowsky-conj}), and the
roots of $Q$ are invariant under conjugation (\Cref{lem:daub-real-poly-roots}).
\end{proof}

\begin{lemma}[Zeros of the Fej\'er--Riesz factor]
\label{lem:daub-fejer-riesz-factor-zeros}
\lean{Polynomial.joukowskyInv_spec_of_mem_aroots, Polynomial.eval_fejerRieszFactor_ne_zero}
\leanok
\uses{def:daub-fejer-riesz-factor}
Let $Q$ be a real polynomial with $Q(r) > 0$ for every $r \in [-1, 1]$. Then every complex root
$t$ of $Q$ satisfies $|a(t)| < 1$ and $a(t)^2 - 2t \, a(t) + 1 = 0$, and $L_0(z) \neq 0$ for
every $z \in \C$ with $|z| \leq 1$.
\end{lemma}

\begin{proof}
\leanok
\uses{lem:daub-real-poly-roots, lem:daub-joukowsky-root}
No root of $Q$ lies in $[-1, 1]$ (\Cref{lem:daub-real-poly-roots}), so the equation for $a(t)$
has a root in the open unit disk (\Cref{lem:daub-joukowsky-root}). For $|z| \leq 1$ each factor
$1 - a(t) z$ of $L_0(z)$ is nonzero, since $|a(t) z| < 1$.
\end{proof}

\begin{lemma}[The Fej\'er--Riesz factor on the unit circle]
\label{lem:daub-fejer-riesz-factor-modulus}
\lean{Polynomial.prod_two_mul_norm_mul_prod_norm_eq}
\leanok
\uses{def:daub-fejer-riesz-factor}
Let $Q$ be a real polynomial with $Q(r) > 0$ for every $r \in [-1, 1]$. Then for every
$\theta \in \R$
\[
  \prod_t 2 |a(t)| \cdot \prod_t |\cos\theta - t| = |L_0(e^{i\theta})|^2 .
\]
\end{lemma}

\begin{proof}
\leanok
\uses{lem:daub-joukowsky-factor, lem:daub-fejer-riesz-factor-zeros, lem:daub-joukowsky-conj,
  lem:daub-real-poly-roots}
Let $z = e^{i\theta}$, so $|z| = 1$ and $\operatorname{Re} z = \cos\theta$. For each root $t$,
\Cref{lem:daub-joukowsky-factor,lem:daub-fejer-riesz-factor-zeros} give
\[
  2 |a(t)| \, |\cos\theta - t| = |1 - a(t) z| \, \bigl|1 - \overline{a(t)} z\bigr| .
\]
Since $\overline{a(t)} = a(\bar t)$ (\Cref{lem:daub-joukowsky-conj}) and the roots are invariant
under conjugation (\Cref{lem:daub-real-poly-roots}),
$\prod_t |1 - \overline{a(t)} z| = \prod_t |1 - a(t) z| = |L_0(z)|$.
\end{proof}

\paragraph{The Fej\'er--Riesz theorem.}
On the unit circle $|L_0|^2$ is $Q(\cos\theta)$ up to a positive constant; normalising the
constant and the sign gives the theorem.

\begin{lemma}[Fej\'er--Riesz up to a constant]
\label{lem:daub-fejer-riesz-constant}
\lean{Polynomial.exists_eval_cos_eq_mul_norm_sq}
\leanok
Let $Q$ be a real polynomial with $Q(r) > 0$ for every $r \in [-1, 1]$. Then there are $K > 0$
and a real polynomial $L$ with $\deg L \leq \deg Q$ and $L(z) \neq 0$ for every $z \in \C$ with
$|z| \leq 1$ such that $Q(\cos\theta) = K \, |L(e^{i\theta})|^2$ for every $\theta \in \R$.
\end{lemma}

\begin{proof}
\leanok
\uses{def:daub-fejer-riesz-factor, lem:daub-fejer-riesz-factor-real,
  lem:daub-fejer-riesz-factor-zeros, lem:daub-fejer-riesz-factor-modulus,
  lem:daub-real-poly-roots, lem:daub-joukowsky-equation}
Let $L$ be $L_0$ read as a real polynomial (\Cref{lem:daub-fejer-riesz-factor-real}); its degree
is at most $\deg Q$ (\Cref{def:daub-fejer-riesz-factor}), and it has no zeros in the closed unit
disk (\Cref{lem:daub-fejer-riesz-factor-zeros}). Since $Q(0) > 0$, $Q \neq 0$; let $c \neq 0$ be its
leading coefficient and $K = |c| / \prod_t 2 |a(t)|$, which is positive since $a(t) \neq 0$
(\Cref{lem:daub-joukowsky-equation}). At $x = \cos\theta \in [-1, 1]$, $Q(x) > 0$, so
\[
  Q(x) = |c| \prod_t |x - t| = K \, |L_0(e^{i\theta})|^2
\]
by \Cref{lem:daub-real-poly-roots} and \Cref{lem:daub-fejer-riesz-factor-modulus}.
\end{proof}

\begin{lemma}[Fixing the sign]
\label{lem:daub-fejer-riesz-sign}
\lean{Polynomial.exists_eval_one_pos_of_aeval_ne_zero}
\leanok
Let $L$ be a real polynomial with $L(z) \neq 0$ for every $z \in \C$ with $|z| \leq 1$. Then
there is a real polynomial $L'$ with $\deg L' = \deg L$ and $L'(1) > 0$, with $L'(z) \neq 0$ for
every $z \in \C$ with $|z| \leq 1$, and with $|L'(z)| = |L(z)|$ for every $z \in \C$.
\end{lemma}

\begin{proof}
\leanok
$L(1) \neq 0$; take $L' = L$ if $L(1) > 0$, and $L' = -L$ otherwise.
\end{proof}

\begin{theorem}[The Fej\'er--Riesz theorem]
\label{thm:daub-fejer-riesz}
\lean{Polynomial.fejer_riesz_outer, Polynomial.fejer_riesz}
\leanok
Let $Q$ be a real polynomial with $Q(t) > 0$ for every $t \in [-1, 1]$. Then there is a real
polynomial $L$ with $\deg L \leq \deg Q$, $L(1) > 0$ and $L(z) \neq 0$ for every $z \in \C$ with
$|z| \leq 1$, such that
\[
  |L(e^{i\theta})|^2 = Q(\cos\theta) \quad\text{for every } \theta \in \R .
\]
\end{theorem}

\begin{proof}
\leanok
\uses{lem:daub-fejer-riesz-constant, lem:daub-fejer-riesz-sign}
Take $K$ and $L$ from \Cref{lem:daub-fejer-riesz-constant}, and apply
\Cref{lem:daub-fejer-riesz-sign} to $\sqrt K L$; then
$|\sqrt K L(e^{i\theta})|^2 = K \, |L(e^{i\theta})|^2 = Q(\cos\theta)$.
\end{proof}

\begin{lemma}[Two readings of the identity]
\label{lem:daub-fejer-riesz-readings}
\lean{Polynomial.forall_norm_aeval_exp_neg_sq_iff, Polynomial.eval_one_sq_of_norm_aeval_exp_sq}
\leanok
Let $Q$ and $L$ be real polynomials. Then $|L(e^{-i\theta})|^2 = Q(\cos\theta)$ for every
$\theta \in \R$ if and only if $|L(e^{i\theta})|^2 = Q(\cos\theta)$ for every $\theta \in \R$,
and then $L(1)^2 = Q(1)$.
\end{lemma}

\begin{proof}
\leanok
Replace $\theta$ by $-\theta$, and use $\cos(-\theta) = \cos\theta$. At $\theta = 0$ the identity
reads $|L(1)|^2 = Q(1)$, and $L(1)$ is real.
\end{proof}

\begin{corollary}[The Fej\'er--Riesz theorem with $e^{-i\theta}$]
\label{cor:daub-fejer-riesz-exp-neg}
\lean{Polynomial.fejer_riesz_exp_neg}
\leanok
Let $Q$ be a real polynomial with $Q(t) > 0$ for every $t \in [-1, 1]$. Then there is a real
polynomial $L$ with $\deg L \leq \deg Q$ and $L(1) > 0$ such that
\[
  |L(e^{-i\theta})|^2 = Q(\cos\theta) \quad\text{for every } \theta \in \R .
\]
\end{corollary}

\begin{proof}
\leanok
\uses{thm:daub-fejer-riesz, lem:daub-fejer-riesz-readings}
\Cref{thm:daub-fejer-riesz} and \Cref{lem:daub-fejer-riesz-readings}.
\end{proof}

With $\theta = 2\pi\xi$ this is the form $|L(e^{-2\pi i \xi})|^2$ of the masks. For example,
$Q(t) = 2 - t = P_2((1 - t)/2)$ has the single root $t = 2$, with $a(2) = 2 - \sqrt 3$, and the
construction gives $L(z) = \frac12 \bigl((1 + \sqrt 3) + (1 - \sqrt 3) z\bigr)$, the polynomial of
the first Daubechies wavelet, $N = 2$.
